\section{The tautological quotient and the universal property of $\mathrm{Gr}(r,d)$}
\label{mk-chap:Picard_GrassmannianQuot}

\providecommand{\ShMod}{\mathrm{SheafOfModules}}

The previous chapter constructed the Grassmannian scheme \(\mathrm{Gr}(r,d)\) over
\(\mathbb{Z}\) by gluing the affine charts \(U^I\)
(\ref{mk-def:gr_affine_chart}) along the transition isomorphisms
(\ref{mk-def:gr_transition}) satisfying the cocycle condition
(\ref{mk-lem:gr_cocycle}), and proved that the structure morphism
\(\mathrm{Gr}(r,d) \to \Spec \mathbb{Z}\) is separated (\ref{mk-lem:gr_separated})
and proper (\ref{mk-lem:gr_proper}). This chapter adds the next layer: the
\emph{tautological rank-\(d\) quotient}
\(u : \mathcal{O}_{\mathrm{Gr}(r,d)}^{\,r} \twoheadrightarrow \mathcal{U}\) of the
trivial rank-\(r\) sheaf, and the \emph{universal property} that makes
\(\mathrm{Gr}(r,d)\) the fine moduli space of rank-\(d\) locally free quotients of
\(\mathcal{O}^{r}\) -- equivalently, the representing scheme of the Quot functor
of the trivial sheaf.

Throughout, \(r \ge d \ge 1\) are the fixed integers of the previous chapter, and
for a size-\(d\) subset \(I \subseteq \{1,\dots,r\}\) we write \(R^I\) for the
chart ring \(\mathbb{Z}[X^I]\) of \(U^I = \Spec R^I\) and \(R^I_J\) for the
localisation along \(P^I_J = \det(X^I_J)\) defining the overlap
\(U^I_J \subseteq U^I\).

\section{Infrastructure: locally free sheaves and module gluing}
\label{mk-sec:grquot_infra}

The constructions of this chapter -- the universal quotient sheaf
\(\mathcal{U}\), the tautological quotient \(u\), and the functor of points -- rest
on two pieces of foundational machinery for sheaves of modules. The first is the
predicate ``locally free of rank \(d\)'' (\ref{mk-def:is_locally_free_of_rank}),
which records that a sheaf of modules is, on the members of some open cover,
isomorphic to the trivial bundle \(\mathcal{O}^{d}\); it is the condition cut out
in the definition of the Grassmannian functor. That predicate is the project-local
\(\mathrm{SheafOfModules.IsLocallyFreeOfRank}\) already set up in the Quot-scheme
chapter (\ref{mk-def:is_locally_free_of_rank}), which this chapter reuses verbatim;
the only foundational machinery genuinely new to this chapter is the second piece:
a gluing primitive that descends a sheaf of modules (and a morphism of such
sheaves) from per-chart data along an open cover, given a transition cocycle
compatible with the cover's own gluing; it is the vehicle that assembles the
per-chart free sheaves \(\mathcal{O}^{d}_{U^I}\) and the per-chart quotients
\(u^I\) into the global \(\mathcal{U}\) and \(u\), and is set up below.

Stating the triple-overlap multiplicativity of the descent datum, however, first
requires a small piece of base-change machinery for the pullback functor on sheaves
of modules: the three transition isomorphisms attached to a triple of charts live a
priori on three different overlaps, and to compare their composite they must be
transported to a single common overlap. The next two blocks supply exactly this
transport -- the pseudofunctorial comparison of iterated pullbacks
(\ref{mk-def:modules_pullbackComp}) and its packaging along the triple-overlap data of a
scheme glue datum (\ref{mk-lem:modules_pullback_basechange_transport}) -- after which the
cocycle condition can be expressed cleanly.

\begin{definition}[Comparison of iterated pullbacks of sheaves of modules]
  \label{mk-def:modules_pullbackComp}
  \lean{AlgebraicGeometry.Scheme.Modules.pullbackComp}
  \mathlibok
  \dcref{custom}

  Let \(a : X \to Y\) and \(b : Y \to Z\) be morphisms of schemes. For a sheaf of
  \(\mathcal{O}_Z\)-modules \(\mathcal{M}\), pulling back first along \(b\) and then
  along \(a\) agrees with pulling back along the composite \(b \circ a\): there is a
  natural isomorphism of sheaves of \(\mathcal{O}_X\)-modules
  \[
    a^{*}\bigl(b^{*}\mathcal{M}\bigr)
      \;\xrightarrow{\ \sim\ }\;
    (b \circ a)^{*}\mathcal{M},
  \]
  natural in \(\mathcal{M}\), expressing that \(\mathcal{M} \mapsto \mathcal{M}\)'s
  pullback is a pseudofunctor on schemes. These comparison isomorphisms satisfy the
  pentagon (associativity) coherence for a triple of composable morphisms. This is the
  pullback half of the pseudofunctoriality of sheaves of modules and is supplied by
  Mathlib.
\end{definition}

\begin{lemma}[Base-change transport of a transition isomorphism to a triple overlap]\label{mk-lem:modules_pullback_basechange_transport}
  \lean{AlgebraicGeometry.Scheme.Modules.pullbackBaseChangeTransport}
  \leanok
  \dcref{custom}

  \uses{mk-def:modules_pullbackComp, mk-def:gr_the_glue_data}
  Let \(D\) be a glue datum of schemes (\ref{mk-def:gr_the_glue_data}) with charts
  \(U_i\), overlaps \(V_{ij}\), overlap immersions \(f_{ij} : V_{ij} \hookrightarrow
  U_i\), transitions \(t_{ij} : V_{ij} \xrightarrow{\sim} V_{ji}\), and triple-overlap
  data \(t'_{ijk}\) with its compatibility identity \(t_{\mathrm{fac}}\). For an index
  triple \((i,j,k)\) write
  \[
    V_{ijk} \;=\; V_{ij} \times_{U_i} V_{ik}
  \]
  for the common triple overlap, with the two projections
  \(p^{ij}_{ijk} : V_{ijk} \to V_{ij}\) and \(p^{ik}_{ijk} : V_{ijk} \to V_{ik}\), and
  let \(t'_{ijk} : V_{ijk} \xrightarrow{\sim} V_{jki}\) be the glue datum's triple
  transition, satisfying
  \[
    t'_{ijk} \,\circ\, p^{jk}_{jki} \;=\; p^{ij}_{ijk} \,\circ\, t_{ij}
    \qquad (t_{\mathrm{fac}}).
  \]
  Given per-chart sheaves of modules \(\mathcal{M}_i\) and a transition isomorphism
  \[
    g_{ij} \;:\; f_{ij}^{*}\,\mathcal{M}_i
      \;\xrightarrow{\ \sim\ }\;
    (t_{ij} \circ f_{ji})^{*}\,\mathcal{M}_j
  \]
  on the overlap \(V_{ij}\) (as in \ref{mk-def:scheme_modules_glue}), the
  \emph{base-change transport} of \(g_{ij}\) to the triple overlap \(V_{ijk}\) is the
  isomorphism of sheaves of \(\mathcal{O}_{V_{ijk}}\)-modules
  \[
    \widehat{g}_{ij}^{\,k}
      \;:\;
    (f_{ij} \circ p^{ij}_{ijk})^{*}\,\mathcal{M}_i
      \;\xrightarrow{\ \sim\ }\;
    (t_{ij} \circ f_{ji} \circ p^{ij}_{ijk})^{*}\,\mathcal{M}_j
  \]
  obtained by pulling \(g_{ij}\) back along \(p^{ij}_{ijk}\) and reassociating the
  iterated pullbacks via \ref{mk-def:modules_pullbackComp}. With the analogous transports
  \(\widehat{g}_{jk}^{\,i}\) and \(\widehat{g}_{ik}^{\,j}\) (the first taken along
  \(t'_{ijk}\), so that \(t_{\mathrm{fac}}\) identifies its source and target sheaves
  with those of \(\widehat{g}_{ij}^{\,k}\) and \(\widehat{g}_{ik}^{\,j}\)), all three
  transported isomorphisms become endomorphism-composable: \(\widehat{g}_{ij}^{\,k}\)
  and \(\widehat{g}_{jk}^{\,i}\) land on a common sheaf, and their composite is
  comparable to \(\widehat{g}_{ik}^{\,j}\) on \(V_{ijk}\).
\end{lemma}

\begin{proof}
  \leanok
  \uses{mk-def:modules_pullbackComp, mk-def:gr_the_glue_data}
  The pullback functor on sheaves of modules is a pseudofunctor: for composable
  scheme morphisms it carries the comparison isomorphisms of
  \ref{mk-def:modules_pullbackComp}, natural in the sheaf and coherent for triple
  composites. Pulling the transition isomorphism \(g_{ij}\) back along the projection
  \(p^{ij}_{ijk} : V_{ijk} \to V_{ij}\) yields an isomorphism between
  \((p^{ij}_{ijk})^{*} f_{ij}^{*}\mathcal{M}_i\) and \((p^{ij}_{ijk})^{*}(t_{ij} \circ
  f_{ji})^{*}\mathcal{M}_j\); applying \ref{mk-def:modules_pullbackComp} to each side
  reassociates these into the single-pullback form
  \((f_{ij} \circ p^{ij}_{ijk})^{*}\mathcal{M}_i\) and
  \((t_{ij} \circ f_{ji} \circ p^{ij}_{ijk})^{*}\mathcal{M}_j\), giving
  \(\widehat{g}_{ij}^{\,k}\). For the \((j,k)\)-transition the same construction is
  applied along the triple transition \(t'_{ijk}\) rather than a bare projection; the
  glue datum's compatibility identity \(t_{\mathrm{fac}}\),
  \(t'_{ijk} \circ p^{jk}_{jki} = p^{ij}_{ijk} \circ t_{ij}\), guarantees that the
  source of \(\widehat{g}_{jk}^{\,i}\) coincides with the target of
  \(\widehat{g}_{ij}^{\,k}\) (both are the pullback of \(\mathcal{M}_j\) along the same
  morphism \(V_{ijk} \to U_j\)), so the two are composable as endomorphism-style
  isomorphisms over \(V_{ijk}\). The reassociations introduced by
  \ref{mk-def:modules_pullbackComp} are coherent (they satisfy the pentagon identity), so
  the composite \(\widehat{g}_{jk}^{\,i} \circ \widehat{g}_{ij}^{\,k}\) and the third
  transport \(\widehat{g}_{ik}^{\,j}\) are isomorphisms between the same pair of
  sheaves on \(V_{ijk}\), and may be compared directly. This is the module-level
  analogue of the scheme-level triangle \(t'_{ijk}\)/\(t_{\mathrm{fac}}\) and is the
  shape in which the multiplicative cocycle condition is stated below.
\end{proof}

\begin{definition}[Gluing a sheaf of modules along an open cover]\label{mk-def:scheme_modules_glue}
  \lean{AlgebraicGeometry.Scheme.Modules.glue}
  \leanok
  \dcref{custom}

  \uses{mk-def:gr_the_glue_data, mk-def:modules_pullbackComp,
    mk-lem:modules_pullback_basechange_transport, mk-lem:gr_glueData_bridges}
  Let \(D\) be a glue datum of schemes (\ref{mk-def:gr_the_glue_data}), with index set \(J\), charts \(U_i\),
  pairwise overlaps \(V_{ij}\), open immersions
  \(f_{ij} : V_{ij} \hookrightarrow U_i\), and transition isomorphisms
  \(t_{ij} : V_{ij} \xrightarrow{\sim} V_{ji}\) satisfying the glue datum's
  cocycle conditions; write \(X\) for the glued scheme and
  \(\iota_i : U_i \hookrightarrow X\) for the canonical open immersions of the
  charts. Suppose given:
  \begin{itemize}
    \item for each \(i \in J\), a sheaf of \(\mathcal{O}_{U_i}\)-modules \(\mathcal{M}_i\) on the chart \(U_i\);
    \item for each pair \((i,j)\), a \emph{transition isomorphism} of sheaves of
      modules on the overlap \(V_{ij}\),
      \[
        g_{ij} \;:\; f_{ij}^{*}\,\mathcal{M}_i
          \;\xrightarrow{\ \sim\ }\;
        (t_{ij} \circ f_{ji})^{*}\,\mathcal{M}_j,
      \]
      comparing the two charts' sheaves after restriction (pullback) to the
      common overlap \(V_{ij}\);
  \end{itemize}
  \emph{subject to the following module cocycle conditions}, compatible with the
  scheme-level cocycle of \(D\) (its triple transition data \(t_{ij}\) and the
  overlap immersions \(f_{ij}\)), which are required as hypotheses on the family
  \((g_{ij})\):
  \begin{enumerate}
    \item[(C1)] \emph{self-identity:} the self-transition is the identity,
      \(g_{ii} = \mathrm{id}\) on \(f_{ii}^{*}\mathcal{M}_i\) (with
      \(V_{ii}\) and \(t_{ii}\) the diagonal data of \(D\));
    \item[(C2)] \emph{triple-overlap multiplicativity:} over each common triple
      overlap \(V_{ijk} = V_{ij} \times_{U_i} V_{ik}\) the base-change transports
      (\ref{mk-lem:modules_pullback_basechange_transport}) of the three transition
      isomorphisms satisfy
      \[
        \widehat{g}_{jk}^{\,i} \circ \widehat{g}_{ij}^{\,k}
          \;=\; \widehat{g}_{ik}^{\,j},
      \]
      where \(\widehat{g}_{ij}^{\,k}\), \(\widehat{g}_{jk}^{\,i}\),
      \(\widehat{g}_{ik}^{\,j}\) are the transports of \(g_{ij}\), \(g_{jk}\),
      \(g_{ik}\) to \(V_{ijk}\) through the scheme-level triple transition
      \(t'_{ijk}\) and its compatibility \(t_{\mathrm{fac}}\), reassociated by the
      pullback comparison \ref{mk-def:modules_pullbackComp} so that all three land on
      the same sheaf -- the multiplicative, \(\mathrm{GL}_d\)-valued cocycle
      condition.
  \end{enumerate}
  Conditions (C1) and (C2) are exactly the descent datum for the Zariski cover
  \(\{\iota_i\}\); it is subject to them that the gluing is well-defined. Out of this
  data the construction produces a glued sheaf of \(\mathcal{O}_X\)-modules
  \(\mathcal{M}\) on the glued scheme \(X\).

  \emph{Construction.} The category \(\ShMod(\mathcal{O}_X)\) of sheaves of
  \(\mathcal{O}_X\)-modules has all limits, and both pushforward along a morphism and
  these limits preserve the sheaf condition; the glued sheaf \(\mathcal{M}\) is therefore
  built directly as a limit -- an \emph{equalizer of pushforwards} -- rather than through a
  hand-built presheaf of compatible families. Concretely, form the two parallel maps of
  sheaves of \(\mathcal{O}_X\)-modules
  \[
    \prod_{i} (\iota_i)_{*}\,\mathcal{M}_i
      \;\;\underset{b}{\overset{a}{\rightrightarrows}}\;\;
    \prod_{i,j} (\iota_i \circ f_{ij})_{*}\,\bigl(f_{ij}^{*}\mathcal{M}_i\bigr),
  \]
  where, on the \((i,j)\)-component, leg \(a\) pushes the \(i\)-th factor forward along the
  overlap immersion \(f_{ij}\) using the unit of the pullback--pushforward adjunction
  together with the pushforward-composition comparison
  \((\iota_i)_{*} \cong (\iota_i \circ f_{ij})_{*}\circ f_{ij}^{*}\) (the dual of
  \ref{mk-def:modules_pullbackComp}), and leg \(b\) pushes the \(j\)-th factor forward and
  additionally transports it across the inverse transition isomorphism \((g_{ij})^{-1}\)
  and the glue datum's compatibility identity, so that both legs land on the common overlap
  sheaf \((\iota_i \circ f_{ij})_{*}(f_{ij}^{*}\mathcal{M}_i)\). The glued sheaf is the
  equalizer
  \[
    \mathcal{M} \;=\; \operatorname{eq}(a,b) \;\hookrightarrow\; \prod_{i} (\iota_i)_{*}\,\mathcal{M}_i,
  \]
  the sub-sheaf of tuples of pushed-forward chart sections whose two overlap restrictions
  agree -- precisely the compatible families, now realised as a categorical limit rather than
  a bespoke presheaf. The cocycle conditions (C1)/(C2) on the family \((g_{ij})\) are
  \emph{not} consumed in forming this equalizer object: the limit exists for any family of
  transition maps. They are instead what make the equalizer descend correctly, and are
  exactly the hypotheses pinned down downstream when the restriction isomorphisms
  \ref{mk-def:gr_modules_glueRestrictionIso} are produced. This specialises effective descent
  of sheaves of modules to the concrete open cover \(\{\iota_i : U_i \hookrightarrow X\}\)
  supplied by a scheme glue datum, and is the form in which the Grassmannian's
  tautological bundle and quotient are built.

  The glued sheaf comes with three further structures, each recorded as its own block
  downstream: the \emph{restriction isomorphisms} \(\rho_i : \iota_i^{*}\mathcal{M}
  \xrightarrow{\sim} \mathcal{M}_i\) compatible with the \(g_{ij}\) over the overlaps
  \ref{mk-def:gr_modules_glueRestrictionIso}; the \emph{characterisation up to unique
  isomorphism} by the restriction property \ref{mk-lem:gr_modules_glue_unique}; and the
  \emph{descent of morphisms} -- a family \(\phi_i : \mathcal{M}_i \to \mathcal{N}_i\)
  commuting with the transition isomorphisms glues to a unique \(\phi : \mathcal{M} \to
  \mathcal{N}\) with \(\rho_i^{\mathcal{N}} \circ \iota_i^{*}\phi = \phi_i \circ
  \rho_i^{\mathcal{M}}\) \ref{mk-def:gr_modules_glueHom}.
\end{definition}

\begin{definition}[Restriction isomorphisms of the glued sheaf]
  \label{mk-def:gr_modules_glueRestrictionIso}
  \lean{}
  \dcref{custom}
  \uses{mk-def:glueRestrictionIso, mk-def:scheme_modules_glue, mk-lem:over_restrict_pullback_iso,
    mk-lem:modules_pullback_basechange_transport, mk-lem:gr_glueData_bridges}
  Let \(\mathcal{M}\) be the glued sheaf of \ref{mk-def:scheme_modules_glue}. For each
  chart index \(i\) there is a \emph{restriction isomorphism} of sheaves of
  \(\mathcal{O}_{U_i}\)-modules
  \[
    \rho_i \;:\; \iota_i^{*}\,\mathcal{M} \;\xrightarrow{\ \sim\ }\; \mathcal{M}_i,
  \]
  compatible with the transition isomorphisms: over each overlap \(V_{ij}\) the two
  restrictions \(\rho_i\) and \(\rho_j\), transported to the common overlap, differ
  exactly by \(g_{ij}\).
\end{definition}
\begin{proof}
  \uses{mk-def:scheme_modules_glue, mk-lem:over_restrict_pullback_iso,
    mk-lem:modules_pullback_basechange_transport, mk-lem:gr_glueData_bridges}
  By the equalizer-of-pushforwards description \ref{mk-def:scheme_modules_glue}, a section of
  \(\mathcal{M}\) over an open of the form \(\iota_i^{-1}V\), restricted to the chart
  \(U_i\), is recovered from and recovers its \(i\)-th component
  \(s_i \in \Gamma(\mathcal{M}_i, \cdot)\) -- the \(i\)-th factor of the product
  \(\prod_i (\iota_i)_{*}\mathcal{M}_i\) into which the equalizer embeds; the
  chart-restriction bridge \ref{mk-lem:over_restrict_pullback_iso} identifies the pullback
  \(\iota_i^{*}\mathcal{M}\) with this restriction, and the \(i\)-th projection
  \((s_\bullet) \mapsto s_i\) is the isomorphism \(\rho_i\) (with inverse extending a chart
  section by the overlap relation).
  The overlap compatibility is the defining relation of the compatible families: over
  \(V_{ij}\) the components \(s_i, s_j\) are related by \(g_{ij}\), which is precisely the
  statement that \(\rho_i\) and \(\rho_j\) differ by \(g_{ij}\) after transport to
  \(V_{ij}\); the triple-overlap coherence needed for this to be consistent is (C2),
  whose transports are aligned by \ref{mk-lem:modules_pullback_basechange_transport} and
  \ref{mk-lem:gr_glueData_bridges}.
\end{proof}


\begin{definition}[Descent of morphisms along the cover]
  \label{mk-def:gr_modules_glueHom}
  \lean{AlgebraicGeometry.Scheme.Modules.glueHom}
  \lean{}
  \dcref{custom}

  \uses{mk-def:scheme_modules_glue, mk-def:gr_modules_glueRestrictionIso,
    mk-lem:sheaf_existsUnique_gluing_mathlib}
  Let \((\mathcal{M}_i, g_{ij})\) and \((\mathcal{N}_i, h_{ij})\) be two glued data
  over the same glue datum \(D\), with glued sheaves \(\mathcal{M}, \mathcal{N}\) and
  restriction isomorphisms \(\rho^{\mathcal{M}}_i, \rho^{\mathcal{N}}_i\). Given a family
  of \(\mathcal{O}_{U_i}\)-linear morphisms \(\phi_i : \mathcal{M}_i \to \mathcal{N}_i\)
  commuting with the transition isomorphisms on every overlap (i.e.\ \(h_{ij}\) intertwines
  the pullbacks of \(\phi_i\) and \(\phi_j\) over \(V_{ij}\)), there is a unique morphism
  of sheaves of \(\mathcal{O}_X\)-modules \(\phi : \mathcal{M} \to \mathcal{N}\) with
  \[
    \rho^{\mathcal{N}}_i \circ \iota_i^{*}\phi \;=\; \phi_i \circ \rho^{\mathcal{M}}_i
    \qquad \text{for every } i.
  \]
\end{definition}
\begin{proof}
  \uses{mk-def:gr_modules_glueRestrictionIso, mk-lem:sheaf_existsUnique_gluing_mathlib}
  On each chart the prescribed identity defines a local morphism
  \((\rho^{\mathcal{N}}_i)^{-1} \circ \phi_i \circ \rho^{\mathcal{M}}_i :
  \iota_i^{*}\mathcal{M} \to \iota_i^{*}\mathcal{N}\). The commutation of the \(\phi_i\)
  with the transition isomorphisms \((g_{ij}, h_{ij})\), together with the overlap
  compatibility of the \(\rho\)-families (\ref{mk-def:gr_modules_glueRestrictionIso}), makes
  these local morphisms agree on every overlap \(V_{ij}\). By the locality and gluing of
  sections (\ref{mk-lem:sheaf_existsUnique_gluing_mathlib}) applied to the internal hom they
  promote to a unique global morphism \(\phi : \mathcal{M} \to \mathcal{N}\) restricting
  to each \(\phi_i\) through the \(\rho_i\) -- the descent of \(\mathrm{Hom}\) along the
  cover.
\end{proof}

\subsection{Pullback of free sheaves along an arbitrary morphism}
\label{mk-sec:grquot_pullback_free}

The functor of points (\ref{mk-def:grassmannian_functor}) acts on morphisms by pullback, so
it needs the pullback of the trivial bundle to be again trivial --
\(\varphi^{*}\mathcal{O}_T^{\,\oplus I} \cong \mathcal{O}_{T'}^{\,\oplus I}\) -- for an
\emph{arbitrary} scheme morphism \(\varphi\). Mathlib's pullback-of-free comparison applies
only when the underlying site functor \(\mathrm{Opens.map}\,\varphi\) is \emph{final}; the
following three blocks remove that restriction for every morphism, and the resulting
pullback machinery preserves rank-\(d\) local freeness.

\begin{lemma}[The opens-pullback functor is final]\label{mk-lem:gr_opensMap_final}
  \lean{AlgebraicGeometry.Scheme.Modules.opensMap_final}
  \leanok
  \dcref{custom}

  For any morphism of schemes \(\varphi : T' \to T\), the inverse-image functor on opens
  \(\mathrm{Opens.map}\,\varphi : \mathrm{Opens}(T) \to \mathrm{Opens}(T')\),
  \(U \mapsto \varphi^{-1}U\), is \emph{final}.
\end{lemma}

\begin{proof}
  \leanok


Finality means that for every open \(V \subseteq T'\) the structured-arrow category of
  opens \(U \subseteq T\) with \(V \le \varphi^{-1}U\) is connected. That category is
  nonempty and has the terminal object \(U = \top\) -- one always has
  \(V \le \varphi^{-1}\top = T'\), and every object admits a unique arrow to \(\top\) --
  and a category with a terminal object is connected.
\end{proof}

\begin{lemma}[Pullback of a free sheaf is free]\label{mk-lem:gr_pullbackFreeIso}
  \lean{AlgebraicGeometry.Scheme.Modules.pullbackFreeIso}
  \leanok
  \dcref{custom}

  \uses{mk-lem:gr_opensMap_final}
  For any scheme morphism \(\varphi : T' \to T\) and index type \(I\), there is a canonical
  isomorphism of sheaves of \(\mathcal{O}_{T'}\)-modules
  \[
    \varphi^{*}\bigl(\mathcal{O}_T^{\,\oplus I}\bigr)
      \;\xrightarrow{\ \sim\ }\;
    \mathcal{O}_{T'}^{\,\oplus I}.
  \]
\end{lemma}

\begin{proof}
  \leanok
  \uses{mk-lem:gr_opensMap_final}
  Mathlib's pullback-of-free comparison produces this isomorphism whenever the site functor
  \(\mathrm{Opens.map}\,\varphi\) is final; that finality is \ref{mk-lem:gr_opensMap_final},
  which holds for every morphism, so the comparison applies unconditionally.
\end{proof}


\section{The tautological quotient on the charts}
\label{mk-sec:grquot_chart}

\emph{Infrastructure: scalar endomorphisms.} Realising a matrix of regular
functions as a morphism of free sheaves of modules has no off-the-shelf primitive
in Mathlib; it is assembled from two small project-local helpers that turn a global
function into a scalar endomorphism of the structure sheaf, after which the matrix
is reconstituted entry-by-entry over a finite biproduct.

\begin{definition}[Global section of the unit module]\label{mk-def:gr_globalUnitSection}
  \lean{AlgebraicGeometry.Grassmannian.globalUnitSection}
  \leanok
  \dcref{custom}

  \uses{mk-lem:moduleUnit_mathlib}
  Let \(X\) be a scheme and \(a \in \Gamma(X, \mathcal{O}_X)\) a global section of
  the structure sheaf. Viewing \(\mathcal{O}_X\) as the unit module
  \(\mathbf{1} \in \ShMod(\mathcal{O}_X)\) (\ref{mk-lem:moduleUnit_mathlib}), the
  section \(a\) determines a global section of \(\mathbf{1}\): on each open
  \(V \subseteq X\) take the restriction \(a|_V \in \Gamma(V, \mathcal{O}_X)\), the
  family \((a|_V)_V\) being compatible with the restriction maps because
  restriction is functorial. This is the section of the unit sheaf attached to a
  global function.
\end{definition}

\begin{definition}[Scalar endomorphism of \(\mathcal{O}_X\)]\label{mk-def:gr_scalarEnd}
  \lean{AlgebraicGeometry.Grassmannian.scalarEnd}
  \leanok
  \dcref{custom}

  \uses{mk-def:gr_globalUnitSection, mk-lem:moduleUnit_mathlib}
  For \(a \in \Gamma(X, \mathcal{O}_X)\) the \emph{scalar endomorphism}
  \(\mathrm{scalarEnd}(a) : \mathcal{O}_X \to \mathcal{O}_X\) is the
  \(\mathcal{O}_X\)-linear endomorphism of the unit module given by multiplication
  by \(a\). It is obtained from the global section of \ref{mk-def:gr_globalUnitSection}
  under the canonical identification of endomorphisms of the unit module
  \(\mathbf{1}\) with its global sections, \(\End(\mathbf{1}) \cong
  \Gamma(X, \mathbf{1})\); concretely, on each open \(V\) it acts as
  multiplication by \(a|_V\). These scalar endomorphisms are the matrix-entry
  building blocks of the chart quotient.
\end{definition}

\begin{lemma}[Scalar endomorphism of the unit is the identity]\label{mk-lem:gr_scalarEnd_one}
  \lean{AlgebraicGeometry.Grassmannian.scalarEnd_one}
  \leanok
  \dcref{custom}

  \uses{mk-def:gr_scalarEnd}
  The scalar endomorphism attached to the unit function \(1 \in
  \Gamma(X, \mathcal{O}_X)\) is the identity of the unit module:
  \(\mathrm{scalarEnd}(1) = \mathrm{id}_{\mathbf{1}}\).
\end{lemma}

\begin{proof}
  \uses{mk-def:gr_scalarEnd}
  Multiplication by \(1\) is the identity on every section, so under the
  identification \(\End(\mathbf{1}) \cong \Gamma(X, \mathbf{1})\) the section \(1\)
  corresponds to \(\mathrm{id}_{\mathbf{1}}\).
\end{proof}

\begin{lemma}[Scalar endomorphism of the zero function is zero]\label{mk-lem:gr_scalarEnd_zero}
  \lean{AlgebraicGeometry.Grassmannian.scalarEnd_zero}
  \leanok
  \dcref{custom}

  \uses{mk-def:gr_scalarEnd}
  The scalar endomorphism attached to the zero function \(0 \in
  \Gamma(X, \mathcal{O}_X)\) is the zero endomorphism:
  \(\mathrm{scalarEnd}(0) = 0\).
\end{lemma}

\begin{proof}
  \uses{mk-def:gr_scalarEnd}
  Multiplication by \(0\) sends every section to \(0\), so the corresponding
  endomorphism of the unit module is the zero map.
\end{proof}

\begin{lemma}[Composition of scalar endomorphisms is multiplication]\label{mk-lem:gr_scalarEnd_comp}
  \lean{AlgebraicGeometry.Grassmannian.scalarEnd_comp}
  \leanok
  \dcref{custom}

  \uses{mk-def:gr_scalarEnd}
  For \(a, b \in \Gamma(X, \mathcal{O}_X)\) the scalar endomorphisms compose by
  multiplication: \(\mathrm{scalarEnd}(a) \circ \mathrm{scalarEnd}(b) =
  \mathrm{scalarEnd}(a\,b)\).
\end{lemma}
\begin{proof}
  \uses{mk-def:gr_scalarEnd}
  Under the identification \(\End(\mathbf{1}) \cong \Gamma(X, \mathbf{1})\) the
  scalar endomorphisms act as multiplication by \(a\) and by \(b\); their composite
  is multiplication by \(a\,b\), because the section ring is commutative and
  restriction maps are ring homomorphisms.
\end{proof}

\begin{lemma}[Scalar endomorphism is additive]\label{mk-lem:gr_scalarEnd_add}
  \lean{AlgebraicGeometry.Grassmannian.scalarEnd_add}
  \leanok
  \dcref{custom}

  \uses{mk-def:gr_scalarEnd}
  For \(a, b \in \Gamma(X, \mathcal{O}_X)\) one has \(\mathrm{scalarEnd}(a + b) =
  \mathrm{scalarEnd}(a) + \mathrm{scalarEnd}(b)\).
\end{lemma}
\begin{proof}
  \uses{mk-def:gr_scalarEnd}
  Multiplication by \(a + b\) equals multiplication by \(a\) plus multiplication by
  \(b\) on every section, since the restriction maps of \(\mathcal{O}_X\) are additive;
  under \(\End(\mathbf{1}) \cong \Gamma(X, \mathbf{1})\) this is the stated additivity.
\end{proof}


\emph{Infrastructure: matrix automorphisms of the free sheaf.} To realise the
\(\mathrm{GL}_d\) bundle transitions one must promote an invertible \(d \times d\)
matrix of global functions to an \emph{automorphism} of the free rank-\(d\) sheaf
\(\mathcal{O}_S^{\,d}\), exactly as the chart quotient \ref{mk-def:gr_chart_quotient}
realises the universal matrix. The next blocks assemble this matrix-to-morphism
functor and record its two structural identities -- that it carries matrix
multiplication to composition and the identity matrix to the identity -- which are
the algebraic heart of the bundle cocycle.

\begin{definition}[Matrix endomorphism of the free sheaf]\label{mk-def:gr_matrixEnd}
  \lean{AlgebraicGeometry.Grassmannian.matrixEnd}
  \leanok
  \dcref{custom}

  \uses{mk-def:gr_scalarEnd}
  Let \(S\) be a scheme, \(d \in \mathbb{N}\), and
  \(M \in \operatorname{Mat}_{d \times d}(\Gamma(S, \mathcal{O}_S))\) a matrix of
  global functions. The \emph{matrix endomorphism} \(\mathrm{matrixEnd}(M)\) is the
  endomorphism of the free rank-\(d\) sheaf \(\mathcal{O}_S^{\,d}\) whose
  \((p,q)\)-component, under the presentation of \(\mathcal{O}_S^{\,d}\) as the
  finite biproduct of \(d\) copies of the unit module \(\mathbf{1}\), is the scalar
  endomorphism \(\mathrm{scalarEnd}(M_{p,q})\) (\ref{mk-def:gr_scalarEnd}). It is
  assembled from these components by the universal property of the rank-\(d\)
  biproduct, exactly as the chart quotient \(u^I\) is assembled
  (\ref{mk-def:gr_chart_quotient}). This is the substitute for the (non-existent)
  Mathlib ``matrix \(\mapsto\) endomorphism of a free sheaf'' primitive.
\end{definition}

\begin{lemma}[Matrix endomorphism carries multiplication to composition]\label{mk-lem:gr_matrixEnd_comp}
  \lean{AlgebraicGeometry.Grassmannian.matrixEnd_comp}
  \leanok
  \dcref{custom}

  \uses{mk-def:gr_matrixEnd, mk-lem:gr_scalarEnd_comp, mk-lem:gr_scalarEnd_add}
  For \(M, N \in \operatorname{Mat}_{d \times d}(\Gamma(S, \mathcal{O}_S))\) the matrix
  endomorphisms compose according to the matrix product, with the order reversed by
  the contravariance of the component indexing:
  \[
    \mathrm{matrixEnd}(M) \,\circ\, \mathrm{matrixEnd}(N)
      \;=\; \mathrm{matrixEnd}(N \cdot M).
  \]
\end{lemma}
\begin{proof}
  \uses{mk-def:gr_matrixEnd, mk-lem:gr_scalarEnd_comp, mk-lem:gr_scalarEnd_add}
  The composition of two biproduct-assembled morphisms is again biproduct-assembled,
  with \((i,q)\)-component the sum over the middle index \(p\) of the composites of
  components, \(\sum_p \mathrm{component}_{i,p} \circ \mathrm{component}_{p,q}\) (the
  categorical matrix product over the biproduct). Each composite of scalar
  endomorphisms is the scalar endomorphism of the product,
  \(\mathrm{scalarEnd}(a)\circ\mathrm{scalarEnd}(b) = \mathrm{scalarEnd}(ab)\)
  (\ref{mk-lem:gr_scalarEnd_comp}), and a finite sum of scalar endomorphisms is the
  scalar endomorphism of the sum (\ref{mk-lem:gr_scalarEnd_add}); hence the \((i,q)\)-component
  of the composite is \(\mathrm{scalarEnd}\bigl(\sum_p M_{q,p} N_{p,i}\bigr) =
  \mathrm{scalarEnd}\bigl((N\cdot M)_{q,i}\bigr)\), which is precisely the
  \((i,q)\)-component of \(\mathrm{matrixEnd}(N\cdot M)\). The reversal of the factor
  order is the contravariance of the column/component indexing.
\end{proof}

\begin{lemma}[Matrix endomorphism of the identity is the identity]\label{mk-lem:gr_matrixEnd_one}
  \lean{AlgebraicGeometry.Grassmannian.matrixEnd_one}
  \leanok
  \dcref{custom}

  \uses{mk-def:gr_matrixEnd, mk-lem:gr_scalarEnd_one, mk-lem:gr_scalarEnd_zero}
  The matrix endomorphism of the identity matrix is the identity morphism of
  \(\mathcal{O}_S^{\,d}\): \(\mathrm{matrixEnd}(1_{d\times d}) =
  \mathrm{id}_{\mathcal{O}_S^{\,d}}\).
\end{lemma}
\begin{proof}
  \uses{mk-def:gr_matrixEnd, mk-lem:gr_scalarEnd_one, mk-lem:gr_scalarEnd_zero}
  The diagonal components are \(\mathrm{scalarEnd}(1) = \mathrm{id}\)
  (\ref{mk-lem:gr_scalarEnd_one}) and the off-diagonal components are
  \(\mathrm{scalarEnd}(0) = 0\) (\ref{mk-lem:gr_scalarEnd_zero}); the morphism with this
  component matrix is the identity of the biproduct.
\end{proof}

\begin{definition}[Free-sheaf automorphism of an invertible matrix]\label{mk-def:gr_matrixToFreeIso}
  \lean{AlgebraicGeometry.Grassmannian.matrixToFreeIso}
  \leanok
  \dcref{custom}

  \uses{mk-def:gr_matrixEnd, mk-lem:gr_matrixEnd_comp, mk-lem:gr_matrixEnd_one}
  Let \(M, N \in \operatorname{Mat}_{d \times d}(\Gamma(S, \mathcal{O}_S))\) be mutually
  inverse, \(M N = 1\) and \(N M = 1\). The \emph{matrix automorphism}
  \(\mathrm{matrixToFreeIso}(M, N)\) is the isomorphism of free sheaves
  \(\mathcal{O}_S^{\,d} \xrightarrow{\sim} \mathcal{O}_S^{\,d}\) with forward map
  \(\mathrm{matrixEnd}(M)\) and inverse map \(\mathrm{matrixEnd}(N)\); the two
  composites are identities by \ref{mk-lem:gr_matrixEnd_comp} (turning the products
  \(NM\) and \(MN\) into the relevant compositions) and \ref{mk-lem:gr_matrixEnd_one}.
  Its forward map is \(\mathrm{matrixEnd}(M)\) by construction
  (\ref{mk-lem:gr_matrixToFreeIso_hom}).
\end{definition}

\begin{lemma}[Forward map of the matrix automorphism]\label{mk-lem:gr_matrixToFreeIso_hom}
  \lean{AlgebraicGeometry.Grassmannian.matrixToFreeIso_hom}
  \leanok
  \dcref{custom}

  \uses{mk-def:gr_matrixToFreeIso, mk-def:gr_matrixEnd}
  The forward map of \(\mathrm{matrixToFreeIso}(M, N)\) is \(\mathrm{matrixEnd}(M)\).
\end{lemma}
\begin{proof}
  \uses{mk-def:gr_matrixToFreeIso, mk-def:gr_matrixEnd}
  Immediate from the definition \ref{mk-def:gr_matrixToFreeIso}, whose forward field is
  literally \(\mathrm{matrixEnd}(M)\).
\end{proof}

\begin{lemma}[Matrix automorphisms compose multiplicatively]\label{mk-lem:gr_matrixToFreeIso_mul}
  \lean{AlgebraicGeometry.Grassmannian.matrixToFreeIso_mul}
  \leanok
  \dcref{custom}

  \uses{mk-def:gr_matrixToFreeIso, mk-lem:gr_matrixToFreeIso_hom, mk-lem:gr_matrixEnd_comp}
  Let \(A, A', B, B'\) be \(d \times d\) matrices of global functions on \(S\) with
  \(A A' = A' A = 1\), \(B B' = B' B = 1\), so that \(\mathrm{matrixToFreeIso}(A, A')\)
  and \(\mathrm{matrixToFreeIso}(B, B')\) are automorphisms of \(\mathcal{O}_S^{\,d}\).
  Their forward maps compose to the forward map of the product automorphism, with the
  order reversed by the component contravariance:
  \[
    \mathrm{matrixToFreeIso}(A, A')_{\mathrm{hom}}
      \,\circ\, \mathrm{matrixToFreeIso}(B, B')_{\mathrm{hom}}
      \;=\; \mathrm{matrixEnd}(B \cdot A).
  \]
  Equivalently, composition of matrix automorphisms realises matrix multiplication of
  the underlying matrices. This is the linkage that turns the matrix-level cocycle
  identity into a composition identity of sheaf-of-module isomorphisms.
\end{lemma}
\begin{proof}
  \uses{mk-def:gr_matrixToFreeIso, mk-lem:gr_matrixToFreeIso_hom, mk-lem:gr_matrixEnd_comp}
  By \ref{mk-lem:gr_matrixToFreeIso_hom} the two forward maps are
  \(\mathrm{matrixEnd}(A)\) and \(\mathrm{matrixEnd}(B)\); their composite is
  \(\mathrm{matrixEnd}(A)\circ\mathrm{matrixEnd}(B) = \mathrm{matrixEnd}(B\cdot A)\) by
  \ref{mk-lem:gr_matrixEnd_comp}.
\end{proof}

\begin{definition}[Chart quotient \(u^I\)]\label{mk-def:gr_chart_quotient}
  \lean{AlgebraicGeometry.Grassmannian.chartQuotientMap}
  \leanok
  \dcref{custom}

  \uses{mk-def:gr_affine_chart, mk-def:gr_universal_matrix, mk-def:gr_scalarEnd}
  \textit{Source: \cite{nitsure-hilbert-quot}, \S 1.}
  Let \(I \subseteq \{1,\dots,r\}\) with \(\#(I) = d\). On the affine chart
  \(U^I = \Spec R^I\) (\ref{mk-def:gr_affine_chart}) define the
  \(\mathcal{O}_{U^I}\)-linear homomorphism of trivial bundles
  \[
    u^I \;:\; \mathcal{O}_{U^I}^{\,r}
      \;\longrightarrow\; \mathcal{O}_{U^I}^{\,d}
  \]
  given by left multiplication by the universal \(d \times r\) matrix \(X^I\)
  (\ref{mk-def:gr_universal_matrix}): on global sections it is the \(R^I\)-linear
  map \((R^I)^{r} \to (R^I)^{d}\), \(v \mapsto X^I v\). Since the \(I\)-th minor
  \(X^I_I\) of \(X^I\) is the identity matrix \(1_{d \times d}\), the composite of
  \(u^I\) with the inclusion of the \(I\)-indexed coordinates
  \(\mathcal{O}_{U^I}^{\,d} \hookrightarrow \mathcal{O}_{U^I}^{\,r}\) is the
  identity; hence \(u^I\) is a split surjection of locally free sheaves, with
  target the free rank-\(d\) sheaf \(\mathcal{O}_{U^I}^{\,d}\).

  \emph{Realisation.} The free sheaves \(\mathcal{O}_{U^I}^{\,r}\) and
  \(\mathcal{O}_{U^I}^{\,d}\) are the finite biproducts (coproducts) of copies of
  the unit module \(\mathbf{1} = \mathcal{O}_{U^I}\), indexed by \(\{1,\dots,r\}\)
  and \(\{1,\dots,d\}\). Each entry \(X^I_{p,q} \in R^I\) of the universal matrix
  (\ref{mk-def:gr_universal_matrix}) is a global function on \(U^I = \Spec R^I\),
  obtained from \(R^I\) through the canonical identification
  \(\Gamma(\Spec R^I, \mathcal{O}) \cong R^I\); the corresponding scalar
  endomorphism \(\mathrm{scalarEnd}(X^I_{p,q}) : \mathbf{1} \to \mathbf{1}\)
  (\ref{mk-def:gr_scalarEnd}) is the \((p,q)\)-component of \(u^I\). Assembling these
  components over the index biproducts -- the universal property of the biproduct
  packaging a matrix of component morphisms into a single morphism of the free
  sheaves -- yields \(u^I\). This entry-by-entry assembly via the finite biproduct
  is the substitute for a (non-existent) Mathlib ``matrix \(\mapsto\) morphism of
  free sheaves'' primitive.
\end{definition}


\begin{lemma}[The chart quotient is an epimorphism]\label{mk-lem:gr_chartQuotientMap_epi}
  \lean{AlgebraicGeometry.Grassmannian.chartQuotientMap_epi}
  \leanok
  \dcref{custom}

  \uses{mk-def:gr_chart_quotient}
  The chart quotient \(u^I : \mathcal{O}_{U^I}^{\,r} \to \mathcal{O}_{U^I}^{\,d}\)
  (\ref{mk-def:gr_chart_quotient}) is an epimorphism of sheaves of modules.
\end{lemma}

\begin{proof}
  \uses{mk-def:gr_chart_quotient}
  We exhibit a section, so that \(u^I\) is a \emph{split} epimorphism and hence an
  epimorphism. Let \(\iota_I : \{1,\dots,d\} \hookrightarrow \{1,\dots,r\}\) be the
  inclusion of the \(I\)-indexed coordinates, \(k \mapsto \iota_I(k)\), the
  \(k\)-th element of \(I\) in increasing order. It induces a morphism of free
  sheaves \(s_I : \mathcal{O}_{U^I}^{\,d} \to \mathcal{O}_{U^I}^{\,r}\), the
  coordinate inclusion sending the \(k\)-th basis section to the
  \(\iota_I(k)\)-th. We claim
  \[
    u^I \circ s_I \;=\; \mathrm{id}_{\mathcal{O}_{U^I}^{\,d}}.
  \]
  Both sides are morphisms out of the free rank-\(d\) sheaf, which is the coproduct
  of \(d\) copies of the unit module; by the universal property of that coproduct
  it suffices to check the equality on each coordinate, i.e.\ to compare component
  matrices. The \((p,k)\)-component of \(u^I \circ s_I\) is the entry
  \(X^I_{p,\iota_I(k)}\) of the universal matrix at the \(I\)-indexed column, which
  is precisely the \(I\)-th minor \(X^I_I\) of \(X^I\). Since that minor is the
  \(d \times d\) identity (the columns of \(X^I\) indexed by \(I\), read through the
  order isomorphism \(\{1,\dots,d\} \cong I\), form the identity block --
  \ref{mk-def:gr_universal_matrix}), one has \(X^I_{p,\iota_I(k)} = \delta_{p,k}\),
  so \(u^I \circ s_I\) has the identity component matrix and equals
  \(\mathrm{id}\). Hence \(s_I\) splits \(u^I\) and \(u^I\) is a split
  epimorphism, in particular an epimorphism.
\end{proof}

\section{The \(\mathrm{GL}_d\) bundle transition cocycle}
\label{mk-sec:grquot_bundle_cocycle}

The universal quotient sheaf \(\mathcal{U}\) (\ref{mk-def:gr_universal_quotient_sheaf}) is
obtained by feeding the gluing primitive \ref{mk-def:scheme_modules_glue} the per-chart free
rank-\(d\) sheaves \(\mathcal{O}_{U^I}^{\,d}\), the bundle transitions
\(g_{I,J} = (X^I_J)^{-1}\), and the two module cocycle hypotheses (C1)/(C2) that primitive
requires. The matrix cocycle for the Cramer inverses \((X^I_J)^{-1}\) is already in hand --
it is the content of the scheme-level cocycle \ref{mk-lem:gr_cocycle} and the minor identities
\ref{mk-lem:gr_universalMinorInv_identities} -- but the gluing primitive consumes its
\emph{module-level} incarnation: an honest isomorphism of sheaves of modules together with
its self-identity and triple-overlap multiplicativity. The three blocks of this section
package exactly that, decoupling the construction of \(\mathcal{U}\) from the remaining
matrix-to-module transport bookkeeping. This is the only net-new infrastructure the
universal quotient requires beyond the gluing primitive itself.


\begin{definition}[The bundle transition \(g_{I,J}\)]\label{mk-def:gr_bundleTransition}
  \lean{AlgebraicGeometry.Grassmannian.bundleTransition}
  \leanok
  \dcref{custom}

  \uses{mk-def:gr_universalMinorInv, mk-def:gr_scalarEnd, mk-lem:gr_pullbackFreeIso,
    mk-lem:gr_minorDet_unit, mk-def:gr_chart_quotient, mk-def:gr_transition}
  Let \(I, J \subseteq \{1,\dots,r\}\) be size-\(d\) subsets and let
  \(f_{IJ} : U^I_J \hookrightarrow U^I\) and \(t_{IJ} : U^I_J \xrightarrow{\sim} U^J_I\) be
  the overlap immersion and transition of the chart cover (\ref{mk-def:gr_transition}). The
  \emph{bundle transition} \(g_{I,J}\) is the isomorphism of sheaves of modules on the
  overlap \(U^I_J\)
  \[
    g_{I,J} \;:\;
      f_{IJ}^{*}\bigl(\mathcal{O}_{U^I}^{\,d}\bigr)
      \;\xrightarrow{\ \sim\ }\;
      (t_{IJ} \circ f_{JI})^{*}\bigl(\mathcal{O}_{U^J}^{\,d}\bigr)
  \]
  induced by the invertible \(d \times d\) matrix
  \(g_{I,J} = (X^I_J)^{-1} \in \mathrm{GL}_d(R^I_J)\) (\ref{mk-def:gr_universalMinorInv}),
  the Cramer inverse of the localised \(J\)-minor, which is invertible because
  \(\det(X^I_J)\) is a unit of \(R^I_J\) (\ref{mk-lem:gr_minorDet_unit}).

  \emph{Realisation.} The transition is built exactly as the chart quotient
  \ref{mk-def:gr_chart_quotient}: each entry \((X^I_J)^{-1}_{p,q} \in R^I_J\) is a global
  function on the overlap, its scalar endomorphism
  \(\mathrm{scalarEnd}\bigl((X^I_J)^{-1}_{p,q}\bigr)\) (\ref{mk-def:gr_scalarEnd}) is the
  \((p,q)\)-component, and the universal property of the rank-\(d\) biproduct assembles
  these into an automorphism of the free sheaf \(\mathcal{O}_{U^I_J}^{\,d}\); the
  source and target pullbacks of free sheaves are identified with this
  \(\mathcal{O}_{U^I_J}^{\,d}\) through the free-pullback comparisons
  \ref{mk-lem:gr_pullbackFreeIso} along \(f_{IJ}\) and \(t_{IJ} \circ f_{JI}\), conjugating
  the matrix automorphism into the displayed isomorphism. Invertibility of the matrix makes
  \(g_{I,J}\) an isomorphism.
\end{definition}

\begin{lemma}[Self-identity of the bundle transition (C1)]\label{mk-lem:gr_bundleCocycle_id}
  \lean{AlgebraicGeometry.Grassmannian.bundleTransition_self}
  \leanok
  \dcref{custom}

  \uses{mk-def:gr_bundleTransition, mk-lem:gr_universalMatrix_submatrix_self,
    mk-lem:gr_universalMinorInv_identities, mk-lem:gr_scalarEnd_one, mk-lem:gr_scalarEnd_zero,
    mk-lem:gr_pullbackFreeIso}
  The self-transition is the identity: \(g_{I,I} = \mathrm{id}\) on
  \(f_{II}^{*}\mathcal{O}_{U^I}^{\,d}\). This is the module-level form of the descent
  condition (C1) of \ref{mk-def:scheme_modules_glue} for the bundle data \((g_{I,J})\).
\end{lemma}
\begin{proof}
  \uses{mk-def:gr_bundleTransition, mk-lem:gr_universalMatrix_submatrix_self,
    mk-lem:gr_universalMinorInv_identities, mk-lem:gr_scalarEnd_one, mk-lem:gr_scalarEnd_zero}
  The \(I\)-th minor of the universal matrix is the identity block, \(X^I_I = 1_{d\times d}\)
  (\ref{mk-lem:gr_universalMatrix_submatrix_self}), so its Cramer inverse is again the
  identity, \((X^I_I)^{-1} = 1_{d\times d}\) (\ref{mk-lem:gr_universalMinorInv_identities}).
  Hence in the biproduct presentation of \ref{mk-def:gr_bundleTransition} the diagonal
  components are \(\mathrm{scalarEnd}(1) = \mathrm{id}\) (\ref{mk-lem:gr_scalarEnd_one}) and the
  off-diagonal components are \(\mathrm{scalarEnd}(0) = 0\) (\ref{mk-lem:gr_scalarEnd_zero}), so
  the assembled automorphism is the identity of \(\mathcal{O}_{U^I_I}^{\,d}\); conjugation by
  the free-pullback comparisons preserves the identity, giving \(g_{I,I} = \mathrm{id}\).
\end{proof}


\begin{lemma}[Cramer-inverse cocycle on the triple overlap]\label{mk-lem:gr_bundleCocycle_matrix}
  \lean{AlgebraicGeometry.Grassmannian.bundleTransition_cocycle_matrix}
  \leanok
  \dcref{custom}

  \uses{mk-def:gr_universalMinorInv, mk-lem:gr_cocycle_imageMatrix_eq,
    mk-lem:gr_universalMinorInv_identities, mk-lem:gr_mul_submatrix_col,
    mk-lem:gr_inv_mul_inv_mul_cancel, mk-lem:gr_cocycle}
  Over the common triple-overlap ring (the localisation carrying the three charts'
  data to \(V_{IJK}\)) the Cramer inverses of the localised minors satisfy the
  multiplicative cocycle identity
  \[
    (X^J_K)^{-1}\,(X^I_J)^{-1} \;=\; (X^I_K)^{-1}.
  \]
  This is the pure matrix-algebra core of (C2), independent of any sheaf data.
\end{lemma}
\begin{proof}
  \uses{mk-def:gr_universalMinorInv, mk-lem:gr_cocycle_imageMatrix_eq,
    mk-lem:gr_universalMinorInv_identities, mk-lem:gr_mul_submatrix_col,
    mk-lem:gr_inv_mul_inv_mul_cancel, mk-lem:gr_cocycle}
  Write all matrices over the common triple-overlap ring. The image matrix of the
  transition \(\theta_{I,J}\) is \(X^J = (X^I_J)^{-1} X^I\), and its \(K\)-minor is
  computed by restricting to the \(K\)-indexed columns; since taking a column
  submatrix commutes with left multiplication (\ref{mk-lem:gr_mul_submatrix_col}),
  \[
    X^J_K \;=\; \bigl((X^I_J)^{-1} X^I\bigr)_K \;=\; (X^I_J)^{-1}\, X^I_K .
  \]
  Both \(X^I_J\) and \(X^I_K\) have unit determinant
  (\ref{mk-lem:gr_universalMinorInv_identities}), so this product of invertibles inverts
  by \((AB)^{-1} = B^{-1} A^{-1}\):
  \[
    (X^J_K)^{-1} \;=\; (X^I_K)^{-1}\,(X^I_J).
  \]
  Multiplying on the right by \((X^I_J)^{-1}\) and cancelling \(X^I_J (X^I_J)^{-1} = 1\)
  via the two-sided Cramer identity (\ref{mk-lem:gr_universalMinorInv_identities}, the
  cancellation packaged in \ref{mk-lem:gr_inv_mul_inv_mul_cancel}) gives
  \((X^J_K)^{-1} (X^I_J)^{-1} = (X^I_K)^{-1}\). The agreement of the three image
  matrices over the triple overlap that licenses these manipulations is the content of
  \ref{mk-lem:gr_cocycle_imageMatrix_eq}, the matrix form of the scheme-level cocycle
  \ref{mk-lem:gr_cocycle}.
\end{proof}

\begin{lemma}[Scalar endomorphism is natural under pullback]\label{mk-lem:gr_scalarEnd_pullback}
  \lean{AlgebraicGeometry.Grassmannian.scalarEnd_pullback}
  \leanok
  \dcref{custom}

  \uses{mk-def:gr_scalarEnd, mk-lem:gr_pullbackFreeIso, mk-def:modules_pullbackComp}
  Let \(p : T \to S\) be a morphism of schemes and \(a \in \Gamma(S, \mathcal{O}_S)\) a
  global function, with comorphism on global sections
  \(p^{\sharp} : \Gamma(S, \mathcal{O}_S) \to \Gamma(T, \mathcal{O}_T)\).
  Write \(q : p^{*}\mathbf{1} \xrightarrow{\sim} \mathbf{1}\) for the unit-pullback comparison
  (\ref{mk-lem:gr_pullbackFreeIso} at a one-element index). Then the pullback of the scalar endomorphism
  \(\mathrm{scalarEnd}(a)\) is conjugate, through \(q\), to the scalar endomorphism of the
  base-changed function:
  \[
    p^{*}\bigl(\mathrm{scalarEnd}(a)\bigr)
      \;=\; q^{-1} \,\circ\, \mathrm{scalarEnd}\bigl(p^{\sharp} a\bigr) \,\circ\, q
  \]
  as endomorphisms of \(p^{*}\mathbf{1}\); equivalently
  \(q \circ p^{*}(\mathrm{scalarEnd}(a)) = \mathrm{scalarEnd}(p^{\sharp} a) \circ q\). This is the
  single-entry naturality atom underlying the matrix-level statement
  \ref{mk-lem:gr_matrixEnd_pullback}.
\end{lemma}
\begin{proof}
  \uses{mk-def:gr_scalarEnd, mk-lem:gr_pullbackFreeIso}
  The scalar endomorphism \(\mathrm{scalarEnd}(a)\) (\ref{mk-def:gr_scalarEnd}) is multiplication
  by the global section \(a\) on the unit module \(\mathbf 1 = \mathcal{O}_S\). Pulling back
  along \(p\) sends multiplication by \(a\) to multiplication by its image, and the comorphism
  \(p^{\sharp}\) is precisely the action of pullback on global functions, so
  \(p^{*}(\mathrm{scalarEnd}(a))\) is multiplication by \(p^{\sharp}a\) on \(p^{*}\mathbf 1\). The
  comparison \(q\) is an isomorphism of unit modules intertwining this multiplication on
  \(p^{*}\mathbf 1\) with multiplication by \(p^{\sharp}a\) on \(\mathbf 1\); conjugating back
  along \(q\) therefore yields the claimed identity. The verification is a single computation on
  the unit module, where multiplication by a function commutes with the natural comparison
  \(q\).
\end{proof}

\begin{lemma}[Matrix endomorphism is natural under pullback]\label{mk-lem:gr_matrixEnd_pullback}
  \lean{AlgebraicGeometry.Grassmannian.matrixEnd_pullback}
  \leanok
  \dcref{custom}

  \uses{mk-def:gr_matrixEnd, mk-def:gr_scalarEnd, mk-lem:gr_scalarEnd_pullback, mk-lem:gr_pullbackFreeIso,
    mk-def:modules_pullbackComp}
  Let \(p : T \to S\) be a morphism of schemes, \(d \in \mathbb{N}\), and
  \(M \in \operatorname{Mat}_{d\times d}(\Gamma(S, \mathcal{O}_S))\). Let
  \(p^{\sharp}M \in \operatorname{Mat}_{d\times d}(\Gamma(T, \mathcal{O}_T))\) be the entrywise
  image of \(M\) under the comorphism \(p^{\sharp}\), and let
  \(Q : p^{*}(\mathcal{O}_S^{\,d}) \xrightarrow{\sim} \mathcal{O}_T^{\,d}\) be the free-pullback
  comparison \ref{mk-lem:gr_pullbackFreeIso} at the rank-\(d\) index set. Then the
  pullback of the matrix endomorphism \(\mathrm{matrixEnd}(M)\) is conjugate, through \(Q\), to
  the matrix endomorphism of the base-changed matrix:
  \[
    p^{*}\bigl(\mathrm{matrixEnd}(M)\bigr)
      \;=\; Q^{-1} \,\circ\, \mathrm{matrixEnd}(p^{\sharp}M) \,\circ\, Q .
  \]
\end{lemma}
\begin{proof}
  \uses{mk-def:gr_matrixEnd, mk-lem:gr_scalarEnd_pullback, mk-lem:gr_pullbackFreeIso,
    mk-def:modules_pullbackComp}
  By \ref{mk-def:gr_matrixEnd}, \(\mathrm{matrixEnd}(M)\) is the biproduct-assembled endomorphism
  of \(\mathcal{O}_S^{\,d} = \bigoplus_{k} \mathbf 1\) whose \((k,\ell)\)-component is
  \(\mathrm{scalarEnd}(M_{k,\ell})\). The pullback functor \(p^{*}\) is a left adjoint, hence
  additive and biproduct-preserving; through the coproduct comparison underlying the
  free-pullback isomorphism \ref{mk-lem:gr_pullbackFreeIso} it carries the biproduct presentation
  of \(\mathcal{O}_S^{\,d}\) to that of \(\mathcal{O}_T^{\,d}\) and the assembled morphism to the
  morphism reassembled from the pulled-back components, the reassociation of the iterated
  pullbacks being \ref{mk-def:modules_pullbackComp}. On each single \((k,\ell)\)-component the
  pulled-back scalar endomorphism is, by \ref{mk-lem:gr_scalarEnd_pullback}, the
  comparison-conjugate of \(\mathrm{scalarEnd}(p^{\sharp}M_{k,\ell}) =
  \mathrm{scalarEnd}((p^{\sharp}M)_{k,\ell})\). Reassembling over the biproduct, the conjugating
  one-element comparisons combine into the single free comparison \(Q\), giving
  \(p^{*}(\mathrm{matrixEnd}(M)) = Q^{-1}\circ\mathrm{matrixEnd}(p^{\sharp}M)\circ Q\).
\end{proof}

\begin{lemma}[Affine global-sections comorphism is the inducing ring homomorphism]\label{mk-lem:gr_baseChange_bridge_gammaSpec}
  \lean{AlgebraicGeometry.Grassmannian.baseChange_bridge_gammaSpec}
  \leanok
  \dcref{custom}

  \uses{mk-def:gr_glued_scheme}
  Let \(\varphi : A \to B\) be a homomorphism of commutative rings and
  \(\Spec\varphi : \Spec B \to \Spec A\) the induced morphism of affine schemes. Under the
  natural identification of the global sections of an affine scheme with its coordinate ring
  -- the counit \(\Gamma(\Spec A, \mathcal O) \xrightarrow{\sim} A\) of the
  \(\Gamma \dashv \Spec\) adjunction, and likewise for \(B\) -- the global-sections base-change
  comorphism
  \[
    (\Spec\varphi)^{\sharp} : \Gamma(\Spec A, \mathcal O) \longrightarrow \Gamma(\Spec B, \mathcal O)
  \]
  is identified with \(\varphi\) itself; that is, the square relating
  \((\Spec\varphi)^{\sharp}\) to \(\varphi\) through the two counit isomorphisms commutes.
\end{lemma}
\begin{proof}

  This is the naturality of the counit isomorphism \(\Gamma \circ \Spec \cong \mathrm{id}\) of
  the global-sections / spectrum adjunction: for any ring homomorphism \(\varphi\) the comorphism
  of \(\Spec\varphi\) on global sections is, after transport through the natural counit
  isomorphisms at \(A\) and \(B\), the homomorphism \(\varphi\). No Grassmannian-specific data
  enters; the statement is the affine-scheme naturality square applied to a single ring map.
\end{proof}

\begin{lemma}[First-projection bridge to \texorpdfstring{$\mathrm{awayInclLeft}$}{awayInclLeft}]\label{mk-lem:gr_baseChange_bridge_left}
  \lean{AlgebraicGeometry.Grassmannian.baseChange_bridge_left}
  \leanok
  \dcref{custom}

  \uses{mk-def:gr_the_glue_data, mk-def:gr_chart_incl, mk-def:gr_away_pullback_iso,
    mk-lem:gr_awayPullbackIso_inv_fst, mk-def:gr_away_incl_left,
    mk-lem:gr_baseChange_bridge_gammaSpec}
  Fix indices \(I, J, K\) with triple overlap \(V_{IJK} = U^I_J \times_{U^I} U^I_K\) of the glue
  datum (\ref{mk-def:gr_the_glue_data}), and let
  \(p^{IJ}_{IJK} : V_{IJK} \to U^I_J\) be the first projection (the pullback projection of the
  chart inclusion \(\mathrm{chartIncl}\), \ref{mk-def:gr_chart_incl}). Through the away-pullback
  identification \(V_{IJK} \cong \Spec R^I[1/(P^I_J P^I_K)]\) (\ref{mk-def:gr_away_pullback_iso})
  the projection \(p^{IJ}_{IJK}\) is the \(\Spec\) morphism of the left away-inclusion
  \(\mathrm{awayInclLeft}(P^I_J, P^I_K)\) (\ref{mk-lem:gr_awayPullbackIso_inv_fst},
  \ref{mk-def:gr_away_incl_left}). Consequently the global-sections base-change map of
  \(p^{IJ}_{IJK}\), transported through the affine global-sections identification, equals the
  ring homomorphism \(\mathrm{awayInclLeft}(P^I_J, P^I_K)\); in particular the entrywise image of
  the Cramer inverse \((X^I_J)^{-1}\) under the geometric base-change agrees with its image under
  \(\mathrm{awayInclLeft}\).
\end{lemma}
\begin{proof}
  \uses{mk-lem:gr_awayPullbackIso_inv_fst, mk-def:gr_away_incl_left,
    mk-lem:gr_baseChange_bridge_gammaSpec}
  By \ref{mk-lem:gr_awayPullbackIso_inv_fst} the first pullback projection, precomposed with the
  away-pullback identification \ref{mk-def:gr_away_pullback_iso}, equals
  \(\Spec(\mathrm{awayInclLeft}(P^I_J, P^I_K))\). The source and overlap charts are affine, so by
  \ref{mk-lem:gr_baseChange_bridge_gammaSpec} the global-sections comorphism of this \(\Spec\)
  morphism is, through the affine global-sections identification, the ring homomorphism
  \(\mathrm{awayInclLeft}(P^I_J, P^I_K)\) (\ref{mk-def:gr_away_incl_left}). Applying the comorphism
  entrywise to \((X^I_J)^{-1}\) gives the asserted agreement.
\end{proof}

\begin{lemma}[Second-projection bridge to \texorpdfstring{$\mathrm{awayInclRight}$}{awayInclRight}]\label{mk-lem:gr_baseChange_bridge_right}
  \lean{AlgebraicGeometry.Grassmannian.baseChange_bridge_right}
  \leanok
  \dcref{custom}

  \uses{mk-def:gr_the_glue_data, mk-def:gr_chart_incl, mk-def:gr_away_pullback_iso,
    mk-lem:gr_awayPullbackIso_inv_snd, mk-def:gr_away_incl_right,
    mk-lem:gr_baseChange_bridge_gammaSpec}
  With the notation of \ref{mk-lem:gr_baseChange_bridge_left}, let
  \(p^{IK}_{IJK} : V_{IJK} \to U^I_K\) be the second projection (the pullback projection of the
  chart inclusion \(\mathrm{chartIncl}\), \ref{mk-def:gr_chart_incl}). Through the away-pullback
  identification \(V_{IJK} \cong \Spec R^I[1/(P^I_J P^I_K)]\) (\ref{mk-def:gr_away_pullback_iso})
  the projection \(p^{IK}_{IJK}\) is the \(\Spec\) morphism of the right away-inclusion
  \(\mathrm{awayInclRight}(P^I_J, P^I_K)\) (\ref{mk-lem:gr_awayPullbackIso_inv_snd},
  \ref{mk-def:gr_away_incl_right}). Consequently the global-sections base-change map of
  \(p^{IK}_{IJK}\), transported through the affine global-sections identification, equals the
  ring homomorphism \(\mathrm{awayInclRight}(P^I_J, P^I_K)\); in particular the entrywise image of
  the Cramer inverse \((X^I_K)^{-1}\) under the geometric base-change agrees with its image under
  \(\mathrm{awayInclRight}\).
\end{lemma}
\begin{proof}
  \uses{mk-lem:gr_awayPullbackIso_inv_snd, mk-def:gr_away_incl_right,
    mk-lem:gr_baseChange_bridge_gammaSpec}
  Identical to \ref{mk-lem:gr_baseChange_bridge_left} with the second leg in place of the first:
  by \ref{mk-lem:gr_awayPullbackIso_inv_snd} the second pullback projection, precomposed with the
  away-pullback identification \ref{mk-def:gr_away_pullback_iso}, equals
  \(\Spec(\mathrm{awayInclRight}(P^I_J, P^I_K))\). The charts being affine,
  \ref{mk-lem:gr_baseChange_bridge_gammaSpec} identifies its global-sections comorphism with
  \(\mathrm{awayInclRight}(P^I_J, P^I_K)\) (\ref{mk-def:gr_away_incl_right}), and the entrywise
  image of \((X^I_K)^{-1}\) follows.
\end{proof}

\begin{lemma}[Triple-transition bridge to \texorpdfstring{$\Theta_{IJ}$}{Theta_IJ}]\label{mk-lem:gr_baseChange_bridge_transition}
  \lean{AlgebraicGeometry.Grassmannian.baseChange_bridge_transition}
  \leanok
  \dcref{custom}

  \uses{mk-def:gr_the_glue_data, mk-def:gr_chart_transition', mk-def:gr_cocycle_theta_ij,
    mk-def:gr_away_pullback_iso, mk-def:gr_away_mul_comm_equiv,
    mk-lem:gr_baseChange_bridge_gammaSpec}
  With the notation of \ref{mk-lem:gr_baseChange_bridge_left}, let
  \(t'_{IJK} : V_{IJK} \to U^J_K \times_{U^J} U^J_I\) be the triple transition of the glue datum
  (\ref{mk-def:gr_chart_transition'}). Through the source and target away-pullback identifications
  (\ref{mk-def:gr_away_pullback_iso}) and the order-swap isomorphism
  (\ref{mk-def:gr_away_mul_comm_equiv}), \(t'_{IJK}\) is the \(\Spec\) morphism of the localised
  triple-overlap transition \(\Theta_{IJ} : S_J \to S_I\) (\ref{mk-def:gr_cocycle_theta_ij},
  composed with the order-swap ring isomorphism). Consequently the global-sections base-change
  map of \(t'_{IJK}\), transported through the affine global-sections identification, equals the
  ring homomorphism \(\Theta_{IJ}\) (up to the order-swap); in particular the entrywise image of
  the Cramer inverse \((X^J_K)^{-1}\) under the geometric base-change agrees with its image under
  \(\Theta_{IJ}\).
\end{lemma}
\begin{proof}
  \uses{mk-def:gr_chart_transition', mk-def:gr_cocycle_theta_ij, mk-def:gr_away_mul_comm_equiv,
    mk-lem:gr_baseChange_bridge_gammaSpec}
  By \ref{mk-def:gr_chart_transition'} the triple transition \(t'_{IJK}\) is the composite of the
  source away-pullback identification, \(\Spec\Theta_{IJ}\) (\ref{mk-def:gr_cocycle_theta_ij}), the
  \(\Spec\) of the order-swap isomorphism (\ref{mk-def:gr_away_mul_comm_equiv}), and the inverse
  target away-pullback identification. The charts and overlaps being affine, applying
  \ref{mk-lem:gr_baseChange_bridge_gammaSpec} to each \(\Spec\) factor identifies the global-sections
  comorphism of \(t'_{IJK}\), through the affine global-sections identifications and the order-swap,
  with the ring homomorphism \(\Theta_{IJ}\). The entrywise image of \((X^J_K)^{-1}\) under this
  comorphism therefore agrees with its image under \(\Theta_{IJ}\).
\end{proof}

\begin{lemma}[Base-change bridge to the localised cocycle ring homomorphisms]\label{mk-lem:gr_baseChange_bridge}
  \lean{AlgebraicGeometry.Grassmannian.baseChange_bridge}
  \leanok
  \dcref{custom}

  \uses{mk-def:gr_transition, mk-def:gr_glued_scheme, mk-def:gr_the_glue_data,
    mk-lem:gr_bundleCocycle_matrix, mk-def:gr_cocycle_theta_ij, mk-def:gr_away_incl_left,
    mk-def:gr_away_incl_right, mk-lem:gr_baseChange_bridge_left,
    mk-lem:gr_baseChange_bridge_right, mk-lem:gr_baseChange_bridge_transition}
  Fix indices \(I, J, K\) and the triple overlap \(V_{IJK}\) of the chart cover
  \(\{U^I\}\) of \(\mathrm{Gr}(r,d)\) (\ref{mk-def:gr_glued_scheme}, \ref{mk-def:gr_the_glue_data}).
  The scheme-pullback base-change map on global sections
  \(\Gamma(U^I_J, \mathcal{O}) \to \Gamma(V_{IJK}, \mathcal{O})\) induced by the first projection
  \(p^{IJ}_{IJK}\), and likewise the maps induced by the second projection and by the triple
  transition \(t'_{IJK}\) (\ref{mk-def:gr_transition}), are identified -- through the natural identification of global sections of an affine scheme
  with its coordinate ring and the fact that the projections and triple transition are the
  \(\Spec\) morphisms of the coordinate-ring homomorphisms by which the Grassmannian glue datum
  (\ref{mk-def:gr_the_glue_data}) is assembled -- with the ring homomorphisms
  \(\mathrm{awayInclLeft}\) (\ref{mk-def:gr_away_incl_left}),
  \(\mathrm{awayInclRight}\) (\ref{mk-def:gr_away_incl_right}) and \(\Theta_{IJ}\)
  (\ref{mk-def:gr_cocycle_theta_ij}) over which the matrix cocycle
  \ref{mk-lem:gr_bundleCocycle_matrix} is stated. Consequently the entrywise image of a Cramer
  inverse \((X^I_J)^{-1}\) under the geometric base-change map agrees with its image under the
  corresponding coordinate-ring homomorphism, so \ref{mk-lem:gr_bundleCocycle_matrix} applies
  to the base-changed matrices.
\end{lemma}
\begin{proof}
  \uses{mk-lem:gr_baseChange_bridge_left, mk-lem:gr_baseChange_bridge_right,
    mk-lem:gr_baseChange_bridge_transition}
  The three identifications are exactly the three projection bridges. The first projection is
  handled by \ref{mk-lem:gr_baseChange_bridge_left}, identifying its global-sections base-change
  map with \(\mathrm{awayInclLeft}\); the second projection by
  \ref{mk-lem:gr_baseChange_bridge_right}, identifying it with \(\mathrm{awayInclRight}\); and the
  triple transition \(t'_{IJK}\) by \ref{mk-lem:gr_baseChange_bridge_transition}, identifying it
  with \(\Theta_{IJ}\). Each bridge already records the agreement of the entrywise image of the
  relevant Cramer inverse under the geometric base-change with its image under the named ring
  homomorphism over which \ref{mk-lem:gr_bundleCocycle_matrix} is stated, so assembling the three
  gives the claim.
\end{proof}

\begin{lemma}[Transport and endpoint alignment of the bundle transitions]\label{mk-lem:gr_bundleCocycle_transport}
  \lean{AlgebraicGeometry.Grassmannian.bundleTransition_cocycle_transport}
  \leanok
  \dcref{custom}

  \uses{mk-def:gr_bundleTransition, mk-lem:gr_bundleCocycle_matrix, mk-lem:gr_matrixToFreeIso_mul,
    mk-lem:gr_matrixEnd_pullback, mk-lem:gr_baseChange_bridge,
    mk-lem:modules_pullback_basechange_transport, mk-lem:gr_glueData_bridges,
    mk-def:modules_pullbackComp, mk-lem:gr_pullbackFreeIso}
  Fix indices \(I, J, K\) with triple overlap \(V_{IJK} = U^I_J \times_{U^I} U^I_K\),
  and let \(\widehat{g}_{IJ}^{\,K}, \widehat{g}_{JK}^{\,I}, \widehat{g}_{IK}^{\,J}\) be
  the base-change transports (\ref{mk-lem:modules_pullback_basechange_transport}) of the
  three bundle transitions to \(V_{IJK}\). Under the free-pullback comparisons
  (\ref{mk-lem:gr_pullbackFreeIso}) and the endpoint bridges
  (\ref{mk-lem:gr_glueData_bridges}), each transported transition is identified with the
  matrix automorphism of the corresponding base-changed Cramer inverse on the common
  free sheaf \(\mathcal{O}_{V_{IJK}}^{\,d}\); granting the matrix identity
  \ref{mk-lem:gr_bundleCocycle_matrix}, the composite satisfies
  \[
    \widehat{g}_{JK}^{\,I} \circ \widehat{g}_{IJ}^{\,K}
      \;=\; \widehat{g}_{IK}^{\,J}
  \]
  as isomorphisms of sheaves of modules on \(V_{IJK}\). This is the substantive
  transport step: the alignment of the three transports' domains and codomains.
\end{lemma}
\begin{proof}
  \uses{mk-def:gr_bundleTransition, mk-lem:gr_bundleCocycle_matrix, mk-lem:gr_matrixToFreeIso_mul,
    mk-lem:gr_matrixEnd_pullback, mk-lem:gr_baseChange_bridge,
    mk-lem:modules_pullback_basechange_transport, mk-lem:gr_glueData_bridges,
    mk-def:modules_pullbackComp, mk-lem:gr_pullbackFreeIso}
  This is pure assembly of the two infrastructure pieces over the endpoint bridges. Each
  bundle transition is, by \ref{mk-def:gr_bundleTransition}, \(\mathrm{matrixToFreeIso}\) of a
  localised Cramer inverse conjugated by the free-pullback comparisons
  \ref{mk-lem:gr_pullbackFreeIso}; its base-change transport
  (\ref{mk-lem:modules_pullback_basechange_transport}) reassociates the iterated pullbacks via
  \ref{mk-def:modules_pullbackComp}, and the three endpoint bridges \ref{mk-lem:gr_glueData_bridges}
  pin the source, middle and target sheaves to a single \(\mathcal{O}_{V_{IJK}}^{\,d}\). By
  \ref{mk-lem:gr_matrixEnd_pullback} each transport then equals \(\mathrm{matrixEnd}\) of the
  base-changed Cramer inverse, and by \ref{mk-lem:gr_baseChange_bridge} those base-changed matrices
  are exactly \((X^I_J)^{-1}, (X^J_K)^{-1}, (X^I_K)^{-1}\) over the common triple-overlap ring of
  \ref{mk-lem:gr_bundleCocycle_matrix}. The linkage \ref{mk-lem:gr_matrixToFreeIso_mul} composes the
  first two to \(\mathrm{matrixEnd}\bigl((X^J_K)^{-1}(X^I_J)^{-1}\bigr)\), and
  by the matrix cocycle identity \ref{mk-lem:gr_bundleCocycle_matrix} the argument equals
  \((X^I_K)^{-1}\), so the composite equals \(\widehat{g}_{IK}^{\,J}\); the bridges make both
  sides isomorphisms between the same pair of sheaves on \(V_{IJK}\).
\end{proof}

\begin{lemma}[Triple-overlap multiplicativity of the bundle transition (C2)]\label{mk-lem:gr_bundleCocycle_mul}
  \lean{AlgebraicGeometry.Grassmannian.bundleTransition_cocycle}
  \leanok
  \dcref{custom}

  \uses{mk-def:gr_bundleTransition, mk-lem:gr_bundleCocycle_matrix, mk-lem:gr_matrixToFreeIso_mul,
    mk-lem:gr_bundleCocycle_transport, mk-lem:modules_pullback_basechange_transport,
    mk-lem:gr_glueData_bridges, mk-def:modules_pullbackComp}
  Over each triple overlap \(V_{IJK} = U^I_J \times_{U^I} U^I_K\) the base-change transports
  (\ref{mk-lem:modules_pullback_basechange_transport}) of the three bundle transitions satisfy
  \[
    \widehat{g}_{JK}^{\,I} \circ \widehat{g}_{IJ}^{\,K}
      \;=\; \widehat{g}_{IK}^{\,J},
  \]
  the module-level form of the descent condition (C2) of \ref{mk-def:scheme_modules_glue} for
  the bundle data \((g_{I,J})\). This is the \(\mathrm{GL}_d\)-valued multiplicative cocycle
  of \(\mathcal{U}\).
\end{lemma}
\begin{proof}
  \uses{mk-lem:gr_bundleCocycle_matrix, mk-lem:gr_matrixToFreeIso_mul,
    mk-lem:gr_bundleCocycle_transport}
  The statement decomposes into three independent steps. First, the matrix-algebra core
  \ref{mk-lem:gr_bundleCocycle_matrix} establishes the Cramer-inverse cocycle
  \((X^J_K)^{-1}(X^I_J)^{-1} = (X^I_K)^{-1}\) over the triple overlap. Second, the linkage
  \ref{mk-lem:gr_matrixToFreeIso_mul} records that composition of matrix automorphisms realises
  matrix multiplication. Third, the transport step \ref{mk-lem:gr_bundleCocycle_transport}
  carries the three transitions to the common overlap \(V_{IJK}\), aligns their source,
  middle, and target sheaves through the base-change transport and the endpoint bridges, and
  -- feeding it the matrix identity through the linkage -- yields exactly the stated equality
  \(\widehat{g}_{JK}^{\,I} \circ \widehat{g}_{IJ}^{\,K} = \widehat{g}_{IK}^{\,J}\).
\end{proof}

\begin{definition}[The universal quotient sheaf \(\mathcal{U}\)]\label{mk-def:gr_universal_quotient_sheaf}
  \lean{AlgebraicGeometry.Grassmannian.universalQuotient}
  \leanok
  \dcref{custom}

  \uses{mk-def:gr_chart_quotient, mk-def:gr_universalMinorInv, mk-def:gr_transition,
    mk-lem:gr_cocycle, mk-def:gr_glued_scheme, mk-def:scheme_modules_glue,
    mk-def:gr_modules_glueRestrictionIso, mk-def:gr_bundleTransition,
    mk-lem:gr_bundleCocycle_id, mk-lem:gr_bundleCocycle_mul,
    mk-lem:gr_glueData_bridges, mk-def:is_locally_free_of_rank}
  \textit{Source: \cite{nitsure-hilbert-quot}, \S 1.}
  Over each overlap \(U^I_J = \Spec R^I_J\) define the bundle-transition matrix
  \[
    g_{I,J} \;=\; (X^I_J)^{-1} \;\in\; \mathrm{GL}_d(R^I_J),
  \]
  the Cramer inverse of the localised \(J\)-minor (\ref{mk-def:gr_universalMinorInv}),
  which is invertible because \(\det(X^I_J)\) is a unit of \(R^I_J\). The matrices
  \((g_{I,J})\) are compatible with the chart-gluing cocycle \((\theta_{I,J})\)
  (\ref{mk-def:gr_transition}, \ref{mk-lem:gr_cocycle}): over a triple overlap they
  satisfy \(g_{I,K} = g_{J,K}\, g_{I,J}\) and \(g_{I,I} = 1\), the multiplicative
  cocycle condition for a vector bundle. Gluing the free rank-\(d\) sheaves
  \(\mathcal{O}_{U^I}^{\,d}\) along the isomorphisms \(g_{I,J}\) over the cover
  \(\{U^I\}\) of \(\mathrm{Gr}(r,d)\) (\ref{mk-def:gr_glued_scheme}) therefore
  produces a locally free \(\mathcal{O}_{\mathrm{Gr}(r,d)}\)-module
  \(\mathcal{U}\) of rank \(d\). Its determinant line bundle \(\det \mathcal{U}\)
  has transition functions \(\det(g_{I,J}) = 1/P^I_J\), the data underlying the
  Pl\"ucker embedding.
\end{definition}

The tautological quotient \(u\) is assembled from the chart quotients \(u^I\) by the
section-level descent \ref{mk-def:gr_modules_glueHom}; its one remaining content is the
\emph{overlap-compatibility} of the family \((u^I)\) with the bundle gluing
\((g_{I,J})\). The next four blocks isolate that compatibility: a rectangular analogue of
the square matrix-endomorphism API (\ref{mk-def:gr_matrixEnd},
\ref{mk-lem:gr_matrixEnd_pullback}, \ref{mk-lem:gr_matrixEnd_comp}) for the
\(r \to d\) chart quotient (\ref{mk-def:gr_matrixEndRect},
\ref{mk-lem:gr_matrixEndRect_pullback}, \ref{mk-lem:gr_matrixEndRect_comp}), the
adjunction transposition that turns the descent condition into a pullback-level identity
(\ref{mk-lem:gr_tautologicalQuotientComponent_transpose}), and the matrix-core assembly
that closes it (\ref{mk-lem:gr_tautologicalQuotient_overlap}).

\subsection{Matrix-calculus toolbox}

The chart-cover lemma (\ref{mk-lem:chartLocus_isOpenCover}) and the overlap compatibility of
\ref{mk-def:grPointOfRankQuotient} both rest on a small calculus of \emph{rectangular} matrix
homomorphisms of free sheaves, generalising the square API
(\ref{mk-def:gr_matrixEnd}, \ref{mk-lem:gr_matrixEnd_comp}, \ref{mk-lem:gr_matrixEnd_one}). The
blocks below collect that toolbox: the bridge to the square case, the identity and
composition laws, injectivity of the matrix presentation, the column-restriction law, the
affine splitting of epimorphisms of free sheaves and the resulting matrix right inverse, the
field-level minor extraction, and the presented-matrix machinery that packages a quotient as
a matrix along a morphism and proves Nitsure's change-of-basis identity. They are
project-bespoke infrastructure.


\begin{definition}[Rectangular matrix homomorphism of free sheaves]\label{mk-def:gr_matrixEndRect}
  \lean{AlgebraicGeometry.Grassmannian.matrixEndRect}
  \leanok
  \dcref{custom}

  \uses{mk-def:gr_scalarEnd}
  Let \(S\) be a scheme, \(d, r \in \mathbb{N}\), and
  \(M \in \operatorname{Mat}_{d \times r}(\Gamma(S, \mathcal{O}_S))\) a rectangular matrix
  of global functions. The \emph{rectangular matrix homomorphism}
  \(\mathrm{matrixEndRect}(M)\) is the morphism of free sheaves
  \[
    \mathrm{matrixEndRect}(M) \;:\; \mathcal{O}_S^{\,r}
      \;\longrightarrow\; \mathcal{O}_S^{\,d}
  \]
  whose \((k,\ell)\)-component -- \(k \in \{1,\dots,d\}\) an output index,
  \(\ell \in \{1,\dots,r\}\) an input index -- under the presentation of
  \(\mathcal{O}_S^{\,r}\) and \(\mathcal{O}_S^{\,d}\) as finite biproducts of copies of the
  unit module \(\mathbf 1\), is the scalar endomorphism \(\mathrm{scalarEnd}(M_{k,\ell})\)
  (\ref{mk-def:gr_scalarEnd}). It is assembled from these components by the universal property
  of the two biproducts, exactly as the square \(\mathrm{matrixEnd}\)
  (\ref{mk-def:gr_matrixEnd}) and the chart quotient (\ref{mk-def:gr_chart_quotient}). When
  \(r = d\) it specialises to \(\mathrm{matrixEnd}\); and the chart quotient \(u^I\)
  (\ref{mk-def:gr_chart_quotient}) is by construction \(\mathrm{matrixEndRect}(X^I)\) of the
  universal \(d \times r\) matrix \(X^I\) (\ref{mk-def:gr_universal_matrix}).
\end{definition}

\begin{lemma}[Rectangular matrix homomorphism is natural under pullback]\label{mk-lem:gr_matrixEndRect_pullback}
  \lean{AlgebraicGeometry.Grassmannian.matrixEndRect_pullback}
  \leanok
  \dcref{custom}

  \uses{mk-def:gr_matrixEndRect, mk-def:gr_scalarEnd, mk-lem:gr_scalarEnd_pullback,
    mk-lem:gr_pullbackFreeIso, mk-def:modules_pullbackComp, mk-lem:gr_matrixEnd_pullback}
  Let \(p : T \to S\) be a morphism of schemes and
  \(M \in \operatorname{Mat}_{d\times r}(\Gamma(S,\mathcal{O}_S))\), with entrywise
  base-changed matrix \(p^{\sharp}M \in \operatorname{Mat}_{d\times r}(\Gamma(T,\mathcal{O}_T))\).
  Write \(Q_r : p^{*}\mathcal{O}_S^{\,r} \xrightarrow{\sim} \mathcal{O}_T^{\,r}\) and
  \(Q_d : p^{*}\mathcal{O}_S^{\,d} \xrightarrow{\sim} \mathcal{O}_T^{\,d}\) for the
  free-pullback comparisons (\ref{mk-lem:gr_pullbackFreeIso}) at the index sets
  \(\{1,\dots,r\}\) and \(\{1,\dots,d\}\). Then the pullback of
  \(\mathrm{matrixEndRect}(M)\) is conjugate, through \(Q_r\) and \(Q_d\), to the
  rectangular homomorphism of the base-changed matrix:
  \[
    p^{*}\bigl(\mathrm{matrixEndRect}(M)\bigr)
      \;=\; Q_d^{-1} \,\circ\, \mathrm{matrixEndRect}(p^{\sharp}M) \,\circ\, Q_r .
  \]
\end{lemma}
\begin{proof}
  \uses{mk-def:gr_matrixEndRect, mk-lem:gr_scalarEnd_pullback, mk-lem:gr_pullbackFreeIso,
    mk-def:modules_pullbackComp, mk-lem:gr_matrixEnd_pullback}
  Identical to the square case \ref{mk-lem:gr_matrixEnd_pullback}, with the source and target
  biproducts now indexed by \(\{1,\dots,r\}\) and \(\{1,\dots,d\}\) respectively. The
  pullback functor \(p^{*}\), a left adjoint, is additive and biproduct-preserving; through
  the coproduct comparisons underlying the free-pullback isomorphisms
  \ref{mk-lem:gr_pullbackFreeIso} it carries the biproduct presentation of
  \(\mathcal{O}_S^{\,r}\) (resp. \(\mathcal{O}_S^{\,d}\)) to that of \(\mathcal{O}_T^{\,r}\)
  (resp. \(\mathcal{O}_T^{\,d}\)), the reassociation of the iterated pullbacks being
  \ref{mk-def:modules_pullbackComp}. On each single \((k,\ell)\)-component the pulled-back
  scalar endomorphism is, by \ref{mk-lem:gr_scalarEnd_pullback}, the comparison-conjugate of
  \(\mathrm{scalarEnd}((p^{\sharp}M)_{k,\ell})\). Reassembling over the two biproducts, the
  conjugating one-element comparisons combine into \(Q_r\) on the source and \(Q_d\) on the
  target, giving \(p^{*}(\mathrm{matrixEndRect}(M)) = Q_d^{-1}\circ
  \mathrm{matrixEndRect}(p^{\sharp}M)\circ Q_r\).
\end{proof}

\begin{lemma}[Square-after-rectangular composition law]\label{mk-lem:gr_matrixEndRect_comp}
  \lean{AlgebraicGeometry.Grassmannian.matrixEndRect_comp}
  \leanok
  \dcref{custom}

  \uses{mk-def:gr_matrixEndRect, mk-def:gr_matrixEnd, mk-lem:gr_scalarEnd_comp, mk-lem:gr_scalarEnd_add,
    mk-lem:gr_matrixEnd_comp}
  For \(M \in \operatorname{Mat}_{d\times r}(\Gamma(S,\mathcal{O}_S))\) and
  \(N \in \operatorname{Mat}_{d\times d}(\Gamma(S,\mathcal{O}_S))\), post-composing the
  rectangular homomorphism with a square one realises the matrix product, with the order
  reversed by the contravariance of the component indexing exactly as in
  \ref{mk-lem:gr_matrixEnd_comp}:
  \[
    \mathrm{matrixEndRect}(M) \,\circ\, \mathrm{matrixEnd}(N)
      \;=\; \mathrm{matrixEndRect}(N \cdot M),
  \]
  the composite \(\mathcal{O}_S^{\,r} \xrightarrow{\mathrm{matrixEndRect}(M)}
  \mathcal{O}_S^{\,d} \xrightarrow{\mathrm{matrixEnd}(N)} \mathcal{O}_S^{\,d}\) (here
  \(N \cdot M\) is the \(d \times r\) product).
\end{lemma}
\begin{proof}
  \uses{mk-def:gr_matrixEndRect, mk-def:gr_matrixEnd, mk-lem:gr_scalarEnd_comp, mk-lem:gr_scalarEnd_add,
    mk-lem:gr_matrixEnd_comp}
  The same biproduct-extensionality computation as \ref{mk-lem:gr_matrixEnd_comp}, the only
  change being that the leftmost (input) biproduct is indexed by \(\{1,\dots,r\}\) rather
  than \(\{1,\dots,d\}\). The composite of two biproduct-assembled morphisms is again
  biproduct-assembled, with \((k,\ell)\)-component the sum over the middle index \(m\) of the
  composites of components, \(\sum_m \mathrm{scalarEnd}(N_{k,m}) \circ
  \mathrm{scalarEnd}(M_{m,\ell})\). Each composite of scalar endomorphisms is the scalar
  endomorphism of the product (\ref{mk-lem:gr_scalarEnd_comp}) and a finite sum of scalar
  endomorphisms is the scalar endomorphism of the sum (\ref{mk-lem:gr_scalarEnd_add}); hence
  that component is \(\mathrm{scalarEnd}(\sum_m N_{k,m} M_{m,\ell}) =
  \mathrm{scalarEnd}((N\cdot M)_{k,\ell})\), the \((k,\ell)\)-component of
  \(\mathrm{matrixEndRect}(N\cdot M)\).
\end{proof}

\begin{lemma}[Adjunction transpose of the chart-overlap condition]\label{mk-lem:gr_tautologicalQuotientComponent_transpose}
  \lean{AlgebraicGeometry.Grassmannian.tautologicalQuotientComponent_transpose}
  \leanok
  \dcref{custom}

  \uses{mk-def:gr_chart_quotient, mk-def:gr_bundleTransition, mk-def:gr_modules_glueHom,
    mk-lem:modules_pullback_basechange_transport, mk-lem:gr_homEquiv_conjugateEquiv_app,
    mk-lem:conjugateEquiv_pullbackComp_inv_mathlib, mk-lem:gr_pullbackFreeIso,
    mk-def:modules_pullbackComp}
  Fix size-\(d\) subsets \(I, J\) and let \(f_{IJ} : U^I_J \hookrightarrow U^I\),
  \(t_{IJ} : U^I_J \xrightarrow{\sim} U^J_I\) be the overlap immersion and transition
  (\ref{mk-def:gr_transition}). The per-chart components \(u^I\) of the tautological quotient
  are the adjoint transposes, along the chart immersions, of the chart quotients
  \ref{mk-def:gr_chart_quotient} (precomposed with the free-pullback comparison). Their
  descent compatibility -- the hypothesis the gluing primitive \ref{mk-def:gr_modules_glueHom}
  consumes, namely that the bundle transition \(g_{I,J}\) (\ref{mk-def:gr_bundleTransition})
  intertwines the pullbacks of \(u^I\) and \(u^J\) over \(U^I_J\) -- is equivalent, under the
  pullback--pushforward adjunction of the overlap immersion, to the pullback-level identity
  of morphisms \(\mathcal{O}_{U^I_J}^{\,r} \to \mathcal{O}_{U^I_J}^{\,d}\)
  \[
    g_{I,J} \,\circ\, f_{IJ}^{*}\,u^I
      \;=\; (t_{IJ} \circ f_{JI})^{*}\,u^J,
  \]
  understood through the free-pullback comparisons (\ref{mk-lem:gr_pullbackFreeIso}) on
  source and target.
\end{lemma}
\begin{proof}
  \uses{mk-lem:gr_homEquiv_conjugateEquiv_app, mk-lem:conjugateEquiv_pullbackComp_inv_mathlib,
    mk-lem:gr_pullbackFreeIso, mk-def:modules_pullbackComp}
  The components are produced by applying the adjunction hom-equivalence to the chart
  quotients; the descent condition of \ref{mk-def:gr_modules_glueHom} is stated on those
  transposes. Transposing back across the adjunction is the naturality of the conjugate
  (mate) construction against the hom-equivalences, \ref{mk-lem:gr_homEquiv_conjugateEquiv_app},
  together with the compatibility of the conjugate with the iterated-pullback comparison
  \ref{mk-lem:conjugateEquiv_pullbackComp_inv_mathlib} reassociating the two pullbacks
  \(f_{IJ}^{*}\) and \((t_{IJ}\circ f_{JI})^{*}\) via \ref{mk-def:modules_pullbackComp}. The
  free-pullback comparisons \ref{mk-lem:gr_pullbackFreeIso} re-present both sides as morphisms
  of the trivial bundles \(\mathcal{O}_{U^I_J}^{\,r} \to \mathcal{O}_{U^I_J}^{\,d}\), under
  which the transposed condition becomes the displayed pullback-level identity.
\end{proof}

\begin{lemma}[Overlap compatibility of the tautological quotient]\label{mk-lem:gr_tautologicalQuotient_overlap}
  \lean{AlgebraicGeometry.Grassmannian.tautologicalQuotient_overlap}
  \leanok
  \dcref{custom}

  \uses{mk-lem:gr_tautologicalQuotientComponent_transpose, mk-lem:gr_matrixEndRect_pullback,
    mk-lem:gr_matrixEndRect_comp, mk-def:gr_bundleTransition, mk-def:gr_chart_quotient,
    mk-def:gr_universalMinorInv, mk-lem:gr_universalMatrix_map_transitionPreMap,
    mk-lem:gr_imageMatrix_map_eq, mk-lem:modules_pullback_basechange_transport,
    mk-lem:gr_baseChange_bridge_left}
  For every pair of size-\(d\) subsets \(I, J\), the chart quotients \(u^I\)
  (\ref{mk-def:gr_chart_quotient}) satisfy the overlap-compatibility condition required by the
  section-level descent \ref{mk-def:gr_modules_glueHom} for the bundle gluing \((g_{I,J})\):
  the pullback-level identity \(g_{I,J} \circ f_{IJ}^{*} u^I = (t_{IJ}\circ f_{JI})^{*} u^J\)
  of \ref{mk-lem:gr_tautologicalQuotientComponent_transpose} holds on each overlap \(U^I_J\).
\end{lemma}
\begin{proof}
  \uses{mk-lem:gr_tautologicalQuotientComponent_transpose, mk-lem:gr_matrixEndRect_pullback,
    mk-lem:gr_matrixEndRect_comp, mk-def:gr_bundleTransition,
    mk-lem:gr_universalMatrix_map_transitionPreMap, mk-lem:gr_imageMatrix_map_eq,
    mk-lem:gr_baseChange_bridge_left}
  By \ref{mk-lem:gr_tautologicalQuotientComponent_transpose} it suffices to prove the
  pullback-level identity \(g_{I,J} \circ f_{IJ}^{*} u^I = (t_{IJ}\circ f_{JI})^{*} u^J\) on
  \(U^I_J\). Both sides are morphisms of trivial bundles \(\mathcal{O}_{U^I_J}^{\,r} \to
  \mathcal{O}_{U^I_J}^{\,d}\); write each as \(\mathrm{matrixEndRect}\) of a \(d \times r\)
  matrix over the overlap ring \(R^I_J\) and compare matrices.

  Since \(u^I = \mathrm{matrixEndRect}(X^I)\) (\ref{mk-def:gr_chart_quotient}), the
  base-change naturality \ref{mk-lem:gr_matrixEndRect_pullback} identifies \(f_{IJ}^{*} u^I\),
  through the free-pullback comparisons, with \(\mathrm{matrixEndRect}\) of the localised
  matrix \(X^I\) over \(R^I_J\) (the global-sections base change is the away-localisation
  \ref{mk-lem:gr_baseChange_bridge_left}). The bundle transition \(g_{I,J}\) is
  \(\mathrm{matrixEnd}\) of the Cramer inverse \((X^I_J)^{-1}\)
  (\ref{mk-def:gr_bundleTransition}, \ref{mk-def:gr_universalMinorInv}); hence by the
  square-after-rectangular composition law \ref{mk-lem:gr_matrixEndRect_comp},
  \[
    g_{I,J} \circ f_{IJ}^{*} u^I
      \;=\; \mathrm{matrixEndRect}\bigl((X^I_J)^{-1} \cdot X^I\bigr).
  \]
  The transition image matrix identity \(X^J = (X^I_J)^{-1} X^I\)
  (\ref{mk-lem:gr_universalMatrix_map_transitionPreMap}, transported across the overlap by
  \ref{mk-lem:gr_imageMatrix_map_eq}) turns the right-hand matrix into \(X^J\). Symmetrically,
  \(\mathrm{matrixEndRect}_J\) of \(X^J\) is, by \ref{mk-lem:gr_matrixEndRect_pullback} along
  \(t_{IJ}\circ f_{JI}\), the comparison form of \((t_{IJ}\circ f_{JI})^{*} u^J\). Therefore
  the two sides agree, which is the required overlap compatibility.
\end{proof}

\begin{definition}[The tautological quotient \(u : \mathcal{O}^{r} \twoheadrightarrow \mathcal{U}\)]\label{mk-def:tautological_quotient}
  \lean{AlgebraicGeometry.Grassmannian.tautologicalQuotient}
  \leanok
  \dcref{custom}

  \uses{mk-def:gr_universal_quotient_sheaf, mk-def:gr_chart_quotient,
    mk-def:gr_glued_scheme, mk-def:gr_universalMinorInv, mk-lem:gr_cocycle,
    mk-def:scheme_modules_glue, mk-def:gr_modules_glueRestrictionIso,
    mk-def:gr_modules_glueHom, mk-lem:gr_tautologicalQuotient_overlap}
  \textit{Source: \cite{nitsure-hilbert-quot}, \S 1.}
  The chart quotients \(u^I : \mathcal{O}_{U^I}^{\,r} \to \mathcal{O}_{U^I}^{\,d}\)
  (\ref{mk-def:gr_chart_quotient}) are compatible with the bundle gluing
  \((g_{I,J})\) of \(\mathcal{U}\) (\ref{mk-def:gr_universal_quotient_sheaf}): over
  each overlap \(U^I_J\) one has the identity of \(d \times r\) matrices
  \[
    g_{I,J}\, X^I \;=\; (X^I_J)^{-1} X^I \;=\; X^J \qquad \text{on } U^I_J,
  \]
  which is precisely the defining formula of the chart transition
  \(\theta_{I,J}\) (\ref{mk-def:gr_transition}, \ref{mk-def:gr_universalMinorInv}). At the
  module level this overlap compatibility -- the equalising condition consumed by the
  section-level descent \ref{mk-def:gr_modules_glueHom} -- is exactly
  \ref{mk-lem:gr_tautologicalQuotient_overlap}, whose proof routes the matrix identity above
  through the rectangular pullback/composition API
  (\ref{mk-lem:gr_matrixEndRect_pullback}, \ref{mk-lem:gr_matrixEndRect_comp}) and the
  adjunction transpose \ref{mk-lem:gr_tautologicalQuotientComponent_transpose}.
  Hence the local maps \(u^I\), transported through the trivialisations of
  \(\mathcal{U}\), agree on overlaps and glue to a single surjective
  \(\mathcal{O}_{\mathrm{Gr}(r,d)}\)-linear homomorphism
  \[
    u \;:\; \mathcal{O}_{\mathrm{Gr}(r,d)}^{\,r}
      \;\twoheadrightarrow\; \mathcal{U},
  \]
  the \emph{tautological} (or \emph{universal}) rank-\(d\) quotient of the trivial
  rank-\(r\) sheaf. It is surjective because each \(u^I\) is (split) surjective and
  surjectivity is local on \(\mathrm{Gr}(r,d)\), and \(\mathcal{U}\) is locally
  free of rank \(d\) by \ref{mk-def:gr_universal_quotient_sheaf}. The pair
  \(\langle \mathcal{U}, u \rangle\) is the family of rank-\(d\) quotients carried
  by \(\mathrm{Gr}(r,d)\) itself.
\end{definition}

\section{The functor of points}
\label{mk-sec:grquot_functor}

The Grassmannian functor sends a scheme \(T\) to the set of \emph{equivalence classes} of
rank-\(d\) quotients of \(\mathcal{O}_T^{\,r}\). The following block packages a single such
quotient, the equivalence relation on them, and the pullback action through which the
functor acts on morphisms; the two coherence blocks after it isolate the comparison
identities that make that action functorial.

\begin{definition}[Rank-\(d\) quotient of \(\mathcal{O}_T^{\,r}\)]\label{mk-def:gr_rankQuotient}
  \lean{AlgebraicGeometry.Grassmannian.RankQuotient}
  \leanok
  \dcref{custom}

  A \emph{rank-\(d\) quotient} of the trivial bundle \(\mathcal{O}_T^{\,r}\) on a scheme
  \(T\) is the datum \(x = \langle \mathcal{F}, q\rangle\) of a sheaf of
  \(\mathcal{O}_T\)-modules \(\mathcal{F}\), locally free of rank \(d\), together with an
  epimorphism \(q : \mathcal{O}_T^{\,r} \twoheadrightarrow \mathcal{F}\). Two such quotients
  \(\langle \mathcal{F}, q\rangle\) and \(\langle \mathcal{F}', q'\rangle\) are
  \emph{equivalent} when there is an isomorphism \(f : \mathcal{F} \xrightarrow{\sim}
  \mathcal{F}'\) with \(f \circ q = q'\) (equivalently, \(\ker q = \ker q'\)); this is the
  equivalence relation whose classes form the value of the Grassmannian functor.
\end{definition}

\begin{definition}[The Grassmannian functor \(\mathrm{Grass}(r,d)\)]\label{mk-def:grassmannian_functor}
  \lean{AlgebraicGeometry.Grassmannian.functor}
  \leanok
  \dcref{custom}

  \uses{mk-def:gr_rankQuotient, mk-lem:gr_pullbackFreeIso, mk-lem:gr_pullbackObjUnitToUnit_id,
    mk-lem:gr_pullbackObjUnitToUnit_comp}
  The \emph{Grassmannian functor} \(\mathrm{Grass}(r,d) : \mathbf{Sch}^{\mathrm{op}} \to
  \mathbf{Set}\) sends a scheme \(T\) to the set of equivalence classes of rank-\(d\)
  quotients \(\langle \mathcal{F}, q\rangle\) of \(\mathcal{O}_T^{\,r}\)
  (\ref{mk-def:gr_rankQuotient}), and a morphism \(\psi : T' \to T\) to the \emph{pullback
  action} \(\langle \mathcal{F}, q\rangle \mapsto \langle \psi^{*}\mathcal{F}, \psi^{*}q\rangle\),
  the source being re-identified with \(\mathcal{O}_{T'}^{\,r}\) through the free-pullback
  isomorphism \(\psi^{*}\mathcal{O}_T^{\,r} \cong \mathcal{O}_{T'}^{\,r}\)
  (\ref{mk-lem:gr_pullbackFreeIso}). Pullback preserves epimorphisms and rank-\(d\) local
  freeness, so the action is well defined and respects the equivalence relation, hence
  descends to classes. The identity and composition functoriality laws follow from the
  pseudofunctor coherences of \(\psi^{*}\), reduced by coproduct extensionality over
  \(\mathcal{O}^{\,r} = \bigoplus_r \mathbf{1}\) to the unit-level identities
  \ref{mk-lem:gr_pullbackObjUnitToUnit_id} and \ref{mk-lem:gr_pullbackObjUnitToUnit_comp}.
\end{definition}


The two blocks below record the comparison identities, between the free-pullback isomorphism
\ref{mk-lem:gr_pullbackFreeIso} and Mathlib's pseudofunctor structure on pullback, that make
the descended pullback action of \ref{mk-def:gr_rankQuotient} functorial.


\begin{lemma}[Mate/conjugate compatibility of the adjunction hom-equivalence]\label{mk-lem:gr_homEquiv_conjugateEquiv_app}
  \lean{AlgebraicGeometry.Scheme.Modules.homEquiv_conjugateEquiv_app}
  \leanok
  \dcref{custom}

  Let \(\mathrm{adj}_1, \mathrm{adj}_2\) be two adjunctions and \(\alpha\) a natural
  transformation between their left adjoints. For a morphism \(f\) the hom-equivalence of
  \(\mathrm{adj}_2\) applied to \(\alpha_c \circ f\) equals the hom-equivalence of
  \(\mathrm{adj}_1\) applied to \(f\), post-composed with the component at \(d\) of the
  conjugate (mate) of \(\alpha\). This is the general naturality of the conjugate
  construction against the two hom-equivalences; it is the reusable engine behind
  \ref{mk-lem:gr_pullbackObjUnitToUnit_comp}.
\end{lemma}

\begin{proof}
  \leanok


Both sides unfold to the same pasting of unit/counit triangles for the two adjunctions;
  the equality is the standard compatibility of the mate of \(\alpha\) with the
  hom-set bijections.
\end{proof}

\begin{lemma}[Free-pullback unit coherence at the identity]\label{mk-lem:gr_pullbackObjUnitToUnit_id}
  \lean{AlgebraicGeometry.Scheme.Modules.pullbackObjUnitToUnit_id}
  \leanok
  \dcref{custom}

  The free-pullback comparison morphism on the unit sheaf at the identity morphism of \(T\)
  agrees with the component, at the unit sheaf, of the pseudofunctor identity isomorphism
  \(\mathrm{id}^{*} \cong \mathrm{id}\) of the pullback of sheaves of modules.
\end{lemma}
\begin{proof}
  \leanok

Transpose both sides under the pullback--pushforward adjunction: the comparison becomes the
  canonical unit-to-pushforward-unit map, and the pseudofunctor identity becomes the conjugate
  (mate) of the identity adjunction, which acts as the identity on the unit section. The two
  transposes therefore coincide, so the original maps agree.
\end{proof}

\begin{lemma}[Free-pullback unit coherence at a composite]\label{mk-lem:gr_pullbackObjUnitToUnit_comp}
  \lean{AlgebraicGeometry.Scheme.Modules.gr_pullbackObjUnitToUnit_comp}
  \leanok
  \dcref{custom}

  \uses{mk-lem:gr_homEquiv_conjugateEquiv_app}
  For composable \(a : T_y \to T_x\) and \(b : T_z \to T_y\), the free-pullback comparison on
  the unit sheaf at the composite \(b \circ a\) factors, through the pseudofunctor composition
  isomorphism \((b\circ a)^{*}\cong b^{*}a^{*}\), as the pulled-back comparison of \(a\)
  followed by the comparison of \(b\).
\end{lemma}
\begin{proof}
  \leanok
  \uses{mk-lem:gr_homEquiv_conjugateEquiv_app}
  Transpose under the composite pullback--pushforward adjunction. By
  \ref{mk-lem:gr_homEquiv_conjugateEquiv_app} the left side collapses to the unit-to-pushforward
  map of \(b\circ a\) post-composed with the conjugate of the composition isomorphism -- the
  inverse pushforward-composition comparison -- while the right side collapses, by naturality
  of the hom-equivalence, to the unit map of \(a\) followed by its pushforward along \(b\).
  Both reduce to the unit-section identity, so the two sides agree.
\end{proof}


Together \ref{mk-lem:gr_pullbackObjUnitToUnit_id} and \ref{mk-lem:gr_pullbackObjUnitToUnit_comp}
upgrade to the corresponding free-sheaf coherences -- the free-pullback isomorphism at the
identity equals the component of the pullback-identity isomorphism, and likewise for the
composite -- by coproduct extensionality over the decomposition
\(\mathcal{O}^{\,\oplus I} = \bigoplus_I \mathbf 1\). These coherences discharge both
functoriality laws (identity and composition) of the Grassmannian
functor \ref{mk-def:grassmannian_functor}.


\section{The universal property}
\label{mk-sec:grquot_universal}

The representability theorem \ref{mk-thm:grassmannian_universal_property} has two halves:
the forward map \(\varphi \mapsto \langle \varphi^{*}\mathcal{U}, \varphi^{*}u\rangle\)
and its inverse, which manufactures a morphism \(T \to \mathrm{Gr}(r,d)\) out of a rank-\(d\)
quotient \(\langle \mathcal{F}, q\rangle\) on \(T\). The forward map needs the tautological
quotient to actually be a quotient -- that \(u\) is an epimorphism
(\ref{mk-lem:tautologicalQuotient_epi}). The inverse is built Zariski-locally and glued: for
each size-\(d\) subset \(I\) one isolates the open locus \(T_I \subseteq T\) where the
\(I\)-indexed columns of \(q\) form a basis of \(\mathcal{F}\), classifies \(T_I \to U^I\)
there, and glues over the cover \(\{T_I\}\). The blocks of this section supply the four
ingredients of that inverse: the iso-locus of a sheaf morphism (\ref{mk-def:isoLocus}) and
the fact that a morphism becomes invertible after restriction to its iso-locus
(\ref{mk-lem:isIso_pullback_isoLocus_map}); the chart locus \(T_I\) (\ref{mk-def:chartLocus})
and the proof that the chart loci cover \(T\) (\ref{mk-lem:chartLocus_isOpenCover}); and the
glued classifying morphism itself (\ref{mk-def:grPointOfRankQuotient}).

\begin{lemma}[The tautological quotient is an epimorphism]\label{mk-lem:tautologicalQuotient_epi}
  \lean{AlgebraicGeometry.Grassmannian.tautologicalQuotient_epi}
  \leanok
  \dcref{custom}

  \uses{mk-lem:gr_chartQuotientMap_epi, mk-def:tautological_quotient, mk-lem:gr_modules_glue_unique}
  The tautological quotient \(u : \mathcal{O}_{\mathrm{Gr}(r,d)}^{\,r}
  \twoheadrightarrow \mathcal{U}\) (\ref{mk-def:tautological_quotient}) is an epimorphism of
  sheaves of \(\mathcal{O}_{\mathrm{Gr}(r,d)}\)-modules.
\end{lemma}
\begin{proof}
  \uses{mk-lem:gr_chartQuotientMap_epi, mk-def:tautological_quotient, mk-lem:gr_modules_glue_unique}
  Being an epimorphism of sheaves of modules is local on the base: a morphism is epi
  precisely when its restriction to every member of an open cover is epi, since the
  cokernel is computed chartwise and a sheaf vanishing on a cover vanishes
  (the joint detection of morphisms across the chart cover is
  \ref{mk-lem:gr_modules_glue_unique}). The charts
  \(\{U^I\}\) cover \(\mathrm{Gr}(r,d)\), and over \(U^I\) the tautological quotient,
  transported through the trivialisation \(\mathcal{U}|_{U^I} \cong \mathcal{O}_{U^I}^{\,d}\)
  of \ref{mk-def:gr_universal_quotient_sheaf}, is the chart quotient \(u^I :
  \mathcal{O}_{U^I}^{\,r} \to \mathcal{O}_{U^I}^{\,d}\) (\ref{mk-def:tautological_quotient}).
  Each \(u^I\) is a split epimorphism: its \(I\)-minor is the identity \(1_{d\times d}\),
  so the coordinate inclusion \(\mathcal{O}_{U^I}^{\,d} \hookrightarrow
  \mathcal{O}_{U^I}^{\,r}\) on the \(I\)-columns is a section of \(u^I\), whence \(u^I\) is
  epi (\ref{mk-lem:gr_chartQuotientMap_epi}). As \(u^I\) is epi on every chart, \(u\) is epi.
\end{proof}

\begin{definition}[Iso-locus of a morphism of module sheaves]\label{mk-def:isoLocus}
  \lean{AlgebraicGeometry.Grassmannian.isoLocus}
  \leanok
  \dcref{custom}

  For a morphism \(\varphi : \mathcal{M} \to \mathcal{N}\) of sheaves of
  \(\mathcal{O}_X\)-modules on a scheme \(X\), the \emph{iso-locus}
  \(\mathrm{Iso}(\varphi) \subseteq X\) is the supremum (union) of all opens
  \(U \subseteq X\) over which the restriction \(\varphi|_U\) is an isomorphism -- the largest
  open on which \(\varphi\) is invertible. A point \(t\) lies in \(\mathrm{Iso}(\varphi)\)
  precisely when some open neighbourhood of \(t\) pulls \(\varphi\) back to an isomorphism.
\end{definition}

\begin{lemma}[A morphism is invertible over its iso-locus]\label{mk-lem:isIso_pullback_isoLocus_map}
  \lean{AlgebraicGeometry.Grassmannian.isIso_pullback_isoLocus_map}
  \leanok
  \dcref{custom}

  \uses{mk-def:isoLocus}
  The restriction of \(\varphi : \mathcal{M} \to \mathcal{N}\) to its iso-locus
  \(\mathrm{Iso}(\varphi)\) (\ref{mk-def:isoLocus}) is an isomorphism of sheaves of modules.
\end{lemma}
\begin{proof}
  \uses{mk-def:isoLocus}
  Being an isomorphism of module sheaves is stalk-local. Each point of the iso-locus has an
  open neighbourhood on which \(\varphi\) pulls back to an isomorphism, and restriction along
  an open immersion preserves stalks; hence the stalk of the underlying abelian-sheaf morphism
  of \(\varphi\) is invertible at every point of the locus. The same stalk comparison for the
  locus inclusion shows \(\varphi|_{\mathrm{Iso}(\varphi)}\) is stalkwise invertible, so it is
  an isomorphism of abelian sheaves by the stalk criterion; the forgetful functor to abelian
  sheaves reflects this back to module sheaves, and the restriction--pullback comparison
  transports it to the pullback functor.
\end{proof}

\begin{definition}[Chart locus \(T_I\) of a rank-\(d\) quotient]\label{mk-def:chartLocus}
  \lean{AlgebraicGeometry.Grassmannian.chartLocus}
  \leanok
  \dcref{custom}

  \uses{mk-def:isoLocus, mk-def:gr_rankQuotient}
  For a rank-\(d\) quotient \(x = \langle \mathcal{F}, q\rangle\) on \(T\)
  (\ref{mk-def:gr_rankQuotient}) and a size-\(d\) subset \(I \subseteq \{1,\dots,r\}\), the
  \emph{chart composite} is \(c_I = q \circ s_I : \mathcal{O}_T^{\,d} \to \mathcal{F}\) -- the
  \(I\)-indexed coordinate inclusion \(s_I\) of free sheaves followed by the quotient map. The
  \emph{chart locus} \(T_I \subseteq T\) is the iso-locus \(\mathrm{Iso}(c_I)\)
  (\ref{mk-def:isoLocus}): the largest open of \(T\) on which \(c_I\) restricts to an
  isomorphism.
\end{definition}

\begin{lemma}[The chart loci cover \(T\)]\label{mk-lem:chartLocus_isOpenCover}
  \lean{AlgebraicGeometry.Grassmannian.chartLocus_isOpenCover}
  \leanok
  \dcref{custom}

  \uses{mk-def:chartLocus, mk-def:gr_rankQuotient}
  For a rank-\(d\) quotient \(x\) on \(T\) the chart loci \(\{T_I\}_{\#I = d}\)
  (\ref{mk-def:chartLocus}) form an open cover of \(T\).
\end{lemma}
\begin{proof}
  \uses{mk-def:chartLocus}
  Fix \(t \in T\). By local freeness of \(\mathcal{F}\), choose an affine open \(W \ni t\) on
  which \(\mathcal{F}|_W \cong \mathcal{O}_W^{\,d}\); the chart composite then becomes a
  \(d\times r\) matrix of sections of \(\mathcal{O}_W\). Epimorphy of \(q\) makes this matrix
  surjective at the residue field \(\kappa(t)\), so some size-\(d\) column minor, indexed by a
  subset \(I\), has nonzero determinant at \(t\). Shrinking \(W\) to the basic open where that
  determinant is a unit, the minor is invertible, hence the chart composite \(c_I\) is an
  isomorphism there; thus \(t \in T_I\). As \(t\) was arbitrary, the \(T_I\) cover \(T\).
\end{proof}


\subsection{Overlap compatibility of the chart morphisms}

On the intersection \(P = T_I \times_T T_J\) of two chart loci both chart composites pull
back to isomorphisms, so the pulled-back quotient is presented by both chart matrices. The
blocks below assemble the transports comparing the two presentations and proving the two
chart morphisms agree into the glued scheme: invertibility along a composite
(\ref{mk-lem:gr_isIso_pullback_map_comp}), the presented matrix along a composite
(\ref{mk-lem:gr_presentedMatrix_comp}) and under an equality of base morphisms
(\ref{mk-lem:gr_presentedMatrix_congr}), the chart matrix as a presented matrix
(\ref{mk-lem:gr_chartMatrix_eq_presentedMatrix}), its own minor
(\ref{mk-lem:gr_presentedMatrix_submatrix_self}), the realisation
\(\mathrm{aeval}(M^I)(X^I) = M^I\) (\ref{mk-lem:gr_universalMatrix_map_presentedMatrix}), the
image-matrix transport (\ref{mk-lem:gr_imageMatrix_map_ringHom}), chart-morphism naturality
(\ref{mk-lem:gr_comp_chartMorphism}), the glue-condition identity
(\ref{mk-lem:gr_chart_point_eq}), and finally the overlap compatibility itself
(\ref{mk-lem:gr_chartMorphism_glue_compat}).

\begin{definition}[Presented matrix \(M^I\) of a quotient along a morphism]\label{mk-def:gr_presentedMatrix}
  \lean{AlgebraicGeometry.Grassmannian.presentedMatrix}
  \leanok
  \dcref{custom}

  \uses{mk-def:gr_rankQuotient, mk-def:gr_matrixEndRect, mk-lem:gr_pullbackFreeIso}
  Let \(x = \langle \mathcal{F}, q\rangle\) be a rank-\(d\) quotient on \(T\) and \(j : P \to
  T\) a morphism along which the \(I\)-th chart composite \(c_I\) pulls back to an isomorphism.
  The \emph{presented matrix} \(M^I = \mathrm{presentedMatrix}(x, j, I)\) is the \(d\times r\)
  matrix of sections of \(\mathcal{O}_P\) whose rectangular-endomorphism realisation
  (\ref{mk-def:gr_matrixEndRect}) is the conjugated pullback
  \(Q_r^{-1}\circ j^{*}q \circ (j^{*}c_I)^{-1}\circ Q_d\), with \(Q_\bullet\) the free-pullback
  comparisons (\ref{mk-lem:gr_pullbackFreeIso}). It generalises the chart matrix over \(T_I\)
  (the case \(j = \iota_{T_I}\)); its \(I\)-minor is the identity.
\end{definition}

\begin{lemma}[Pullback-invertibility along a composite]\label{mk-lem:gr_isIso_pullback_map_comp}
  \lean{AlgebraicGeometry.Grassmannian.isIso_pullback_map_comp}
  \leanok
  \dcref{custom}

  If \(a^{*}c\) is an isomorphism of module sheaves, then so is \((p\circ a)^{*}c\) for any
  \(p\), via the pseudofunctor comparison \((p\circ a)^{*}\cong p^{*}a^{*}\).
\end{lemma}
\begin{proof}

  The functor \(p^{*}\) carries the isomorphism \(a^{*}c\) to an isomorphism; transporting
  along the composition comparison identifies \(p^{*}(a^{*}c)\) with \((p\circ a)^{*}c\),
  which is therefore an isomorphism.
\end{proof}

\begin{lemma}[Presented matrix along a composite]\label{mk-lem:gr_presentedMatrix_comp}
  \lean{AlgebraicGeometry.Grassmannian.presentedMatrix_comp}
  \leanok
  \dcref{custom}

  \uses{mk-def:gr_presentedMatrix, mk-def:gr_matrixEndRect}
  For \(p : W \to V\) and \(a : V \to T\), the matrix presenting \(x\) along \(p\circ a\) is
  the entrywise \(\Gamma\)-restriction along \(p\) of the matrix presenting \(x\) along \(a\):
  \(\mathrm{presentedMatrix}(x, p\circ a, I) = p^{\sharp}\bigl(\mathrm{presentedMatrix}(x,
  a, I)\bigr)\).
\end{lemma}
\begin{proof}
  \uses{mk-def:gr_presentedMatrix, mk-def:gr_matrixEndRect}
  Both sides are compared through their rectangular-endomorphism realisations
  (\ref{mk-def:gr_matrixEndRect}), which are injective on matrices. The realisation of the left
  side is the conjugated pullback along \(p\circ a\); folding in the composition comparison
  \((p\circ a)^{*}\cong p^{*}a^{*}\) rewrites it as \(p^{*}\) applied to the realisation of the
  right side, conjugated by the free-pullback comparisons, which is exactly the realisation of
  the entrywise restriction \(p^{\sharp}M^I\).
\end{proof}

\begin{lemma}[Presented matrix depends only on the base morphism]\label{mk-lem:gr_presentedMatrix_congr}
  \lean{AlgebraicGeometry.Grassmannian.presentedMatrix_congr}
  \leanok
  \dcref{custom}

  \uses{mk-def:gr_presentedMatrix}
  If \(j = j'\) as morphisms \(P \to T\), then \(\mathrm{presentedMatrix}(x, j, I) =
  \mathrm{presentedMatrix}(x, j', I)\); the invertibility hypotheses are propositionally
  irrelevant.
\end{lemma}
\begin{proof}
  \uses{mk-def:gr_presentedMatrix}
  Substituting \(j' = j\), the two presented matrices coincide by definition.
\end{proof}

\begin{lemma}[Chart matrix as a presented matrix]\label{mk-lem:gr_chartMatrix_eq_presentedMatrix}
  \lean{AlgebraicGeometry.Grassmannian.chartMatrix_eq_presentedMatrix}
  \leanok
  \dcref{custom}

  \uses{mk-def:gr_presentedMatrix, mk-def:chartLocus}
  The chart matrix \(M^I\) of \(x\) over the chart locus \(T_I\) is the presented matrix of
  \(x\) along the locus inclusion: \(\mathrm{chartMatrix}(x, I) =
  \mathrm{presentedMatrix}(x, \iota_{T_I}, I)\).
\end{lemma}
\begin{proof}
  \uses{mk-def:gr_presentedMatrix, mk-def:chartLocus}
  Immediate from the definitions: the chart matrix is, entrywise, the presented matrix along
  the inclusion of the chart locus.
\end{proof}

\begin{lemma}[The own \(I\)-minor of \(M^I\) is the identity]\label{mk-lem:gr_presentedMatrix_submatrix_self}
  \lean{AlgebraicGeometry.Grassmannian.presentedMatrix_submatrix_self}
  \leanok
  \dcref{custom}

  \uses{mk-def:gr_presentedMatrix, mk-def:gr_matrixEndRect}
  The \(I\)-indexed column minor of \(\mathrm{presentedMatrix}(x, j, I)\) is the \(d\times d\)
  identity matrix.
\end{lemma}
\begin{proof}
  \uses{mk-def:gr_presentedMatrix, mk-def:gr_matrixEndRect}
  Through the rectangular-endomorphism realisation, the \(I\)-minor presents the inverted
  chart composite against itself, i.e. \((j^{*}c_I)^{-1}\circ (j^{*}c_I) = \mathrm{id}\);
  matching realisations and using injectivity of the realisation gives the identity matrix.
\end{proof}

\begin{lemma}[The classifying map realises \(X^I\) as \(M^I\)]\label{mk-lem:gr_universalMatrix_map_presentedMatrix}
  \lean{AlgebraicGeometry.Grassmannian.universalMatrix_map_presentedMatrix}
  \leanok
  \dcref{custom}

  \uses{mk-def:gr_presentedMatrix, mk-lem:gr_presentedMatrix_submatrix_self}
  Applying the ring map \(\mathrm{aeval}(M^I)\) -- the substitution sending each free variable
  \(x^I_{p,q}\) to the corresponding entry of the presented matrix -- entrywise to the
  universal matrix \(X^I\) returns \(M^I\): \(\mathrm{aeval}(M^I)(X^I) = M^I\).
\end{lemma}
\begin{proof}
  \uses{mk-def:gr_presentedMatrix, mk-lem:gr_presentedMatrix_submatrix_self}
  On the free entries \(x^I_{p,q}\) (columns outside \(I\)) the substitution returns the
  matching entry of \(M^I\) by \(\mathrm{aeval}(x^I_{p,q}) = M^I_{p,q}\). On the \(I\)-columns
  both \(X^I_I\) and \(M^I_I\) are the identity (\ref{mk-lem:gr_presentedMatrix_submatrix_self}),
  so the two agree there as well.
\end{proof}

\begin{lemma}[Transport of the image matrix along the structure map]\label{mk-lem:gr_imageMatrix_map_ringHom}
  \lean{AlgebraicGeometry.Grassmannian.imageMatrix_map_ringHom}
  \leanok
  \dcref{custom}

  Let \(I, J\) be size-\(d\) subsets and \(\mathrm{incl} : R^I_J \to D\) a ring map into a
  commutative ring \(D\) whose restriction along \(R^I \to R^I_J\) is \(g : R^I \to D\). Then
  the image matrix \((X^I_J)^{-1}X^I\) pushed forward along \(\mathrm{incl}\) equals
  \((Y_J)^{-1}Y\), where \(Y\) is the universal matrix \(X^I\) base-changed along \(g\) and
  \(Y_J\) its \(J\)-minor.
\end{lemma}
\begin{proof}

  Expand the image matrix as the product of the nonsingular inverse of the \(J\)-minor with
  \(X^I\). A ring map commutes with matrix multiplication and, since the minor determinant is
  a unit, with the nonsingular inverse. Rewriting both factors along \(\mathrm{incl}\) and
  using \(\mathrm{incl}\circ(R^I\to R^I_J) = g\) yields \((Y_J)^{-1}Y\).
\end{proof}

\begin{lemma}[Naturality of the chart morphism]\label{mk-lem:gr_comp_chartMorphism}
  \lean{AlgebraicGeometry.Grassmannian.comp_chartMorphism}
  \leanok
  \dcref{custom}

  \uses{mk-def:gr_presentedMatrix, mk-lem:gr_presentedMatrix_comp,
    mk-lem:gr_chartMatrix_eq_presentedMatrix, mk-def:chartLocus}
  Precomposing the chart morphism \(\varphi_I : T_I \to U^I\) with any \(p : W \to T_I\) is
  the morphism \(W \to U^I\) classified by the matrix presenting \(x\) along
  \(p\circ\iota_{T_I}\); i.e. \(p \circ \varphi_I\) is the \(\Gamma\)--\(\Spec\) classifier of
  \(\mathrm{aeval}\bigl(\mathrm{presentedMatrix}(x, p\circ\iota_{T_I}, I)\bigr)\).
\end{lemma}
\begin{proof}
  \uses{mk-lem:gr_presentedMatrix_comp, mk-lem:gr_chartMatrix_eq_presentedMatrix}
  By naturality of \(\Gamma\)--\(\Spec\), \(p\circ\varphi_I\) is classified by the \(p\)-image
  of the ring map classifying \(\varphi_I\), i.e. by \(\mathrm{aeval}(p^{\sharp}M^I)\). By
  \ref{mk-lem:gr_chartMatrix_eq_presentedMatrix} and \ref{mk-lem:gr_presentedMatrix_comp} one has
  \(p^{\sharp}M^I = \mathrm{presentedMatrix}(x, p\circ\iota_{T_I}, I)\), giving the claim.
\end{proof}

\begin{lemma}[Chart-point identity: the glue condition in classifying form]\label{mk-lem:gr_chart_point_eq}
  \lean{AlgebraicGeometry.Grassmannian.chart_point_eq}
  \leanok
  \dcref{custom}

  Let \(K\) be a commutative ring and \(f : R^I \to K\) a ring map under which the minor
  \(P^I_J\) is a unit. Then the transported point
  \(\mathrm{lift}(f)\circ\tilde\theta_{I,J}\) presents the same \(K\)-point through chart
  \(J\) as \(f\) does through chart \(I\); that is,
  \(\Spec(\mathrm{lift}(f)\circ\tilde\theta_{I,J})\circ\iota_J = \Spec(f)\circ\iota_I\) into
  the glued scheme \(\mathrm{Gr}(r,d)\).
\end{lemma}
\begin{proof}

  Factor the \(J\)-classifier through the localization away from \(P^I_J\): the transition map
  \(\tilde\theta_{I,J}\) composed with the chart inclusion is the chart-transition immersion,
  so the glue condition \(\theta_{I,J}\circ\iota_{JI} = \iota_{IJ}\) of the glue data
  transports \(\Spec(\mathrm{lift}(f))\) on chart \(J\) to \(\Spec(f)\) on chart \(I\). The
  universal property of localization (that \(\mathrm{lift}(f)\) restricts to \(f\)) closes the
  identity.
\end{proof}

\begin{lemma}[Overlap compatibility of the chart morphisms]\label{mk-lem:gr_chartMorphism_glue_compat}
  \lean{AlgebraicGeometry.Grassmannian.chartMorphism_glue_compat}
  \leanok
  \dcref{custom}

  \uses{mk-def:chartLocus, mk-lem:gr_isIso_pullback_map_comp, mk-def:gr_presentedMatrix,
    mk-lem:gr_presentedMatrix_submatrix_self, mk-lem:gr_universalMatrix_map_presentedMatrix,
    mk-lem:gr_imageMatrix_map_ringHom, mk-lem:gr_chart_point_eq}
  On the intersection \(P = T_I \times_T T_J\) of two chart loci, the two composite chart
  morphisms agree into the glued scheme:
  \(\mathrm{pr}_1 \circ (\varphi_I \circ \iota_I) = \mathrm{pr}_2 \circ (\varphi_J \circ
  \iota_J)\).
\end{lemma}
\begin{proof}
  \uses{mk-def:gr_presentedMatrix, mk-lem:gr_presentedMatrix_submatrix_self,
    mk-lem:gr_universalMatrix_map_presentedMatrix, mk-lem:gr_imageMatrix_map_ringHom,
    mk-lem:gr_chart_point_eq, mk-lem:gr_isIso_pullback_map_comp}
  On \(P\) both chart composites \(c_I, c_J\) pull back to isomorphisms
  (\ref{mk-lem:gr_isIso_pullback_map_comp}), so the pulled-back quotient is presented by both
  matrices \(M^I_P\) and \(M^J_P\). The \(J\)-minor of \(M^I_P\) is invertible -- it presents
  \(c_J\) against \(c_I^{-1}\) (\ref{mk-lem:gr_presentedMatrix_submatrix_self}) -- so the
  \(I\)-classifier sends \(P^I_J\) to a unit
  (\ref{mk-lem:gr_universalMatrix_map_presentedMatrix}). The change-of-basis identity
  \(M^J_P = (M^I_{P,J})^{-1} M^I_P\), via \ref{mk-lem:gr_imageMatrix_map_ringHom}, identifies the
  localization-transport \(\mathrm{lift}(\mathrm{aeval}\,M^I_P)\circ\tilde\theta_{I,J}\) of the
  \(I\)-classifier with the \(J\)-classifier. Hence by the chart-point identity
  \ref{mk-lem:gr_chart_point_eq} both composites present the same point into the glued scheme,
  so they agree.
\end{proof}

\begin{definition}[The classifying morphism \(T \to \mathrm{Gr}(r,d)\) of a quotient]\label{mk-def:grPointOfRankQuotient}
  \lean{AlgebraicGeometry.Grassmannian.grPointOfRankQuotient}
  \leanok
  \dcref{custom}

  \uses{mk-def:chartLocus, mk-lem:chartLocus_isOpenCover, mk-lem:gr_chartMorphism_glue_compat}
  For a rank-\(d\) quotient \(x = \langle \mathcal{F}, q\rangle\) on \(T\), the
  \emph{classifying morphism} \(\mathrm{gr}(x) : T \to \mathrm{Gr}(r,d)\) is obtained by
  gluing, over the open cover \(\{T_I\}\) of chart loci (\ref{mk-lem:chartLocus_isOpenCover}),
  the composites \(T_I \xrightarrow{\varphi_I} U^I \hookrightarrow \mathrm{Gr}(r,d)\) of the
  chart morphisms with the glue-data immersions; the overlap compatibility making the gluing
  legitimate is \ref{mk-lem:gr_chartMorphism_glue_compat}. It is the inverse of the forward map
  of the universal property.
\end{definition}


The two blocks below carry the chart data of a quotient through a base change
\(\psi : T' \to T\); they are the bridge underlying the two inverse laws of the
representability theorem.

\begin{lemma}[Presented matrix of a pulled-back quotient]\label{mk-lem:gr_presentedMatrix_rqPullback}
  \lean{AlgebraicGeometry.Grassmannian.presentedMatrix_rqPullback}
  \leanok
  \dcref{custom}

  \uses{mk-def:gr_presentedMatrix, mk-def:gr_rankQuotient, mk-def:gr_matrixEndRect}
  \textit{Source: \cite{nitsure-hilbert-quot}, Construction of Grassmannian.}
  For \(\psi : T' \to T\), a rank-\(d\) quotient \(y\) on \(T\), and a morphism
  \(j : W \to T'\), presenting the pulled-back quotient \(\psi^{*}y\) along \(j\) coincides
  with presenting \(y\) along the composite \(j\circ\psi\):
  \[
    \mathrm{presentedMatrix}(\psi^{*}y, j, I)
      \;=\; \mathrm{presentedMatrix}(y, j\circ\psi, I).
  \]
  This is the transport engine of the inverse laws of the universal property: the universal
  matrix \(X^I\) classifying a point of \(\mathrm{Gr}(r,d)\) is, after base change, given by
  the same ring homomorphism out of \(\mathbb{Z}[X^I]\) regardless of whether one pulls back
  first and then presents, or composes the base morphisms.
\end{lemma}
\begin{proof}
  \uses{mk-def:gr_presentedMatrix, mk-def:gr_rankQuotient, mk-def:gr_matrixEndRect}
  Compare the two matrices through their rectangular-endomorphism realisations
  (\ref{mk-def:gr_matrixEndRect}), which are injective. By definition of the pullback action
  (\ref{mk-def:gr_rankQuotient}), \(\psi^{*}y\) re-presents the source through the free-pullback
  comparison, so its pullback along \(j\) reads \(j^{*}(\psi^{*}y).q = j^{*}P_r^{-1}\circ
  j^{*}\psi^{*}q\) and \(j^{*}c_I(\psi^{*}y) = j^{*}P_d^{-1}\circ j^{*}\psi^{*}c_I(y)\), with
  \(P_\bullet\) the comparisons for \(\psi\). Fold the pseudofunctor composition comparison
  \((j\circ\psi)^{*}\cong j^{*}\psi^{*}\) into both the quotient map and the inverted chart
  composite, using the naturality of its unit and counit: the comparison square moves the
  conjugating isomorphisms across \(j^{*}\psi^{*}q\) and \((j^{*}\psi^{*}c_I)^{-1}\), and the
  composite free-pullback comparisons \(P_\bullet, Q_\bullet\) collapse into the single
  comparison for \(j\circ\psi\). What remains is exactly the conjugated pullback of \(y\)
  along \(j\circ\psi\), i.e. the realisation of \(\mathrm{presentedMatrix}(y, j\circ\psi, I)\).
  As the realisations agree, so do the matrices.
\end{proof}

\begin{lemma}[Left inverse law: the classifier recovers \(\psi\)]\label{mk-lem:gr_grPointOfRankQuotient_rqPullback_tautological}
  \lean{AlgebraicGeometry.Grassmannian.grPointOfRankQuotient_rqPullback_tautological}
  \leanok
  \dcref{custom}

  \uses{mk-def:grPointOfRankQuotient, mk-def:tautological_quotient, mk-def:gr_rankQuotient,
    mk-lem:gr_presentedMatrix_rqPullback, mk-lem:gr_presentedMatrix_comp,
    mk-lem:gr_presentedMatrix_congr, mk-lem:gr_comp_chartMorphism, mk-lem:isIso_pullback_isoLocus_map,
    mk-lem:chartLocus_isOpenCover}
  \textit{Source: \cite{nitsure-hilbert-quot}, Construction of Grassmannian (properness).}
  For any morphism \(\psi : T \to \mathrm{Gr}(r,d)\), classifying the pulled-back tautological
  quotient recovers \(\psi\):
  \[
    \mathrm{gr}\bigl(\psi^{*}\langle \mathcal{U}, u\rangle\bigr) \;=\; \psi .
  \]
\end{lemma}
\begin{proof}
  \uses{mk-def:grPointOfRankQuotient, mk-lem:gr_presentedMatrix_rqPullback,
    mk-lem:gr_presentedMatrix_comp, mk-lem:gr_presentedMatrix_congr, mk-lem:gr_comp_chartMorphism,
    mk-lem:isIso_pullback_isoLocus_map}
  The preimages \(\{\psi^{-1}(\iota_I(U^I))\}\) of the chart-immersion ranges cover \(T\), so
  it suffices to check the equality after restriction to each member \(W = \psi^{-1}
  (\iota_I(U^I))\). On \(W\) the composite \(\iota_W\circ\psi\) factors through the chart
  immersion \(\iota_I\) as \(\rho\circ\iota_I\) for a unique \(\rho : W \to U^I\), and \(W\)
  lies inside the chart locus \(T_I\) of \(\psi^{*}\langle\mathcal{U},u\rangle\) (the
  tautological chart composite is invertible over \(\iota_I(U^I)\), and
  \ref{mk-lem:isIso_pullback_isoLocus_map} transports this), so all relevant chart composites
  pull back to isomorphisms. The matrix presenting \(\psi^{*}\langle\mathcal{U},u\rangle\)
  along \(\iota_W\) equals, by \ref{mk-lem:gr_presentedMatrix_rqPullback} and
  \ref{mk-lem:gr_presentedMatrix_congr}, the matrix presenting the tautological pair along
  \(\iota_W\circ\psi = \rho\circ\iota_I\), which by \ref{mk-lem:gr_presentedMatrix_comp} is the
  \(\rho\)-image of the universal matrix \(X^I\) (the tautological presentation over \(U^I\)).
  Hence by chart-morphism naturality \ref{mk-lem:gr_comp_chartMorphism} the glued morphism
  \(\mathrm{gr}(\psi^{*}\langle\mathcal{U},u\rangle)\) restricted to \(W\) is the classifier of
  that matrix, namely \(\rho\circ\iota_I = \iota_W\circ\psi\). As the members \(W\) cover
  \(T\), the two morphisms agree.
\end{proof}

\begin{theorem}[\(\mathrm{Gr}(r,d)\) represents the Grassmannian functor]\label{mk-thm:grassmannian_universal_property}
  \lean{AlgebraicGeometry.Grassmannian.represents}
  \leanok
  \dcref{custom}

  \uses{mk-def:grassmannian_functor, mk-def:tautological_quotient, mk-def:gr_rankQuotient,
    mk-def:grPointOfRankQuotient, mk-lem:gr_grPointOfRankQuotient_rqPullback_tautological}
  \textit{Source: \cite{nitsure-hilbert-quot}, Exercises (2).}
  The scheme \(\mathrm{Gr}(r,d)\), together with the tautological quotient
  \(\langle \mathcal{U}, u\rangle\), represents the Grassmannian functor
  \(\mathrm{Grass}(r,d)\) (\ref{mk-def:grassmannian_functor}): for every scheme \(T\) there is a
  bijection, natural in \(T\),
  \[
    \mathrm{Hom}(T, \mathrm{Gr}(r,d)) \;\xrightarrow{\ \sim\ }\; \mathrm{Grass}(r,d)(T),
    \qquad \psi \longmapsto \langle \psi^{*}\mathcal{U}, \psi^{*}u\rangle .
  \]
\end{theorem}
\begin{proof}
  \uses{mk-def:grPointOfRankQuotient, mk-lem:gr_grPointOfRankQuotient_rqPullback_tautological,
    mk-lem:tautologicalQuotient_epi}
  The forward map sends \(\psi\) to the pullback \(\psi^{*}\langle\mathcal{U},u\rangle\) of the
  tautological pair; naturality in \(T\) is the pseudofunctoriality of the pullback action
  (\ref{mk-def:grassmannian_functor}) evaluated at the tautological point. The inverse is the
  chart-by-chart classifier \(\mathrm{gr}\) of \ref{mk-def:grPointOfRankQuotient}. The two are
  mutually inverse: the left inverse law
  \(\mathrm{gr}(\psi^{*}\langle\mathcal{U},u\rangle) = \psi\) is
  \ref{mk-lem:gr_grPointOfRankQuotient_rqPullback_tautological}, checked over the cover of
  chart-immersion preimages; the right inverse law
  \(\langle\mathrm{gr}(x)^{*}\mathcal{U}, \mathrm{gr}(x)^{*}u\rangle \sim x\) holds because over
  each chart locus of \(x\) both quotients are presented by the same matrix, so the two
  epimorphism descents are mutually inverse, witnessing the equivalence. Surjectivity of \(u\)
  (\ref{mk-lem:tautologicalQuotient_epi}) makes the tautological pair an honest quotient, so the
  forward map indeed lands in \(\mathrm{Grass}(r,d)(T)\).
\end{proof}
